\chapter{文本分类、匹配与组织}
\label{chap:nlp-classification-matching-organization}

从通知中抽出活动、日期和报名条件以后，还不能直接回答另一组问题：这份通知应分到哪些栏目？两份通知是不是在说同一件事？数千份尚无目录的材料可以怎样分组？第~\ref{chap:nlp-information-extraction}章关注原文中的对象、关系与记录，本章转而考察文本本身的属性、文本之间的关系，以及文本集合的组织结构。处理对象改变了，结果的含义和验收方式也随之改变。

沿用本部分的教学通知：活动编号为\code{ML-0418}的“机器学习读书会”于2025年4月15日发布更正：“机器学习读书会原定于4月17日举行，现改至4月18日（周五）14:00，地点为海棠楼201。校内师生可直接参加，校外来宾须在4月16日18:00前提交报名信息。”这份材料既有时间安排，又有报名说明；删去报名条件的短句仍与它主题相关，却未必是等义复述；把它分进“报名”文档簇，也不意味着每句话都讨论报名。本章其余通知、评论、向量、概率和计数同样是教学构造，不是模型实测结果。

表~\ref{tab:nlp-classification-task-meaning}把这些区别放在同一材料上。前卷的分类器、编码器和聚类算法提供计算部件；本章要补齐的是标签怎样定义、监督怎样形成、分数怎样恢复为可使用结果，以及什么反例能够推翻一次看似合理的判断。

\begin{table*}[htbp]
  \centering
  \begin{tabularx}{\linewidth}{@{}>{\raggedright\arraybackslash}p{16mm}
    >{\raggedright\arraybackslash}p{22mm}>{\raggedright\arraybackslash}X
    >{\raggedright\arraybackslash}p{23mm}>{\raggedright\arraybackslash}X
    >{\raggedright\arraybackslash}p{22mm}@{}}
    \toprule
    任务 & 输入对象 & 输出含义 & 监督依据 & 通知中的反例 & 评价口径 \\
    \midrule
    多标签分类 & 一份通知 & 同时包含哪些内容 & 栏目标签规范 & 有报名词不等于只谈报名 & 逐标签及宏、微F1 \\
    情感 & 评论及评价对象 & 对对象的极性或强度 & 指定粒度的意见 & 抱怨改期不等于反对活动 & 按对象、极性评价 \\
    立场 & 文本与指定主张 & 是否支持该主张 & 目标条件标注 & 语气消极仍可支持改期 & 按目标和标签评价 \\
    相似度 & 两段文字 & 意义接近的程度 & 连续参考评分 & 日期不同仍可能得高分 & Pearson、Spearman \\
    复述 & 两段文字 & 约定范围内是否等义 & 等义句对标注 & 删掉报名截止条件 & F1、条件反例 \\
    推断 & 前提与假设 & 前提支持、排斥或未定 & 有方向的三分类 & 有时间不推出有主讲人 & 三分类与方向测试 \\
    相关性 & 查询与候选 & 对当前需求是否有用 & 查询条件相关性 & 讲读书会却不含新时间 & Recall@k、排序指标 \\
    对齐 & 两组文本单元 & 哪些单元相互对应 & 平行材料或对应标注 & 译文把一句拆成两句 & 覆盖率与配对正确性 \\
    主题 & 文档集合 & 共现结构及混合权重 & 无标签词语观测 & 主题比例不是栏目概率 & 连贯性及路线专属指标 \\
    聚类 & 文档集合 & 在给定距离下的组别 & 表示和聚类准则 & 同簇通知可有不同日期 & 稳定性、轮廓及业务核验 \\
    \bottomrule
  \end{tabularx}
  \caption{属性、关系与组织结果的区别。各行围绕同一活动材料改变输入和判断问题；共享分类头或向量表示，不会使这些标签相互等价。F1的统计单位、相关性的候选集合和组织结果的解释条件都必须随任务记录。}
  \label{tab:nlp-classification-task-meaning}
\end{table*}

\FloatBarrier
\section{文本分类与属性判定}
\label{sec:nlp-classification-attributes}

给一段文字贴标签之前，应先问标签在描述什么。“报名说明”描述内容，“通知”描述体裁，“提醒”描述用途；“不满”描述情绪，“支持改期”描述对指定方案的立场。它们可以同时成立。把所有名称放进同一个互斥类别表，会把本来可以并存的属性强迫成竞争关系。分类单位也须固定：一句表达不满的话，不一定代表整篇通知的语气，更不能代表通知转发者的态度。

\subsection{内容类别与用途判定}
\label{sec:nlp-classification-content}

\term{单标签分类}（single-label classification）要求一个样本只选一个类别；\term{多标签分类}（multi-label classification）允许多个属性同时为真。层次分类则增加类别之间的包含关系，例如“活动通知”之下再分“时间安排”和“报名说明”。这里的包含是标注规范的一部分，不是从词义近似自动推出来的。业务还应约定“其他”“未知”和“证据不足”：它们分别可以表示已知的剩余类别、超出类别表的输入，以及当前文本不足以判定，不能任意合并。

设训练样本为$(x_i,\bm y_i)$，$x_i$是文档，$\bm y_i\in\{0,1\}^{K}$记录$K$个标签。文档表示$\bm h_i\in\mathbb R^d$经过$\bm z_i=\bm W\bm h_i+\bm b$得到标签得分。互斥类别使用
\begin{equation}
  \label{eq:nlp-classification-softmax}
  p(y=k\mid x)=\operatorname{softmax}(\bm z)_k
  =\frac{\exp(z_k)}{\sum_{\ell=1}^{K}\exp(z_\ell)}.
\end{equation}
分母遍历全部候选类别，因此概率之和为1；提高某一类的得分会改变其他类的概率。多标签改用逐标签的
\begin{equation}
  \label{eq:nlp-classification-sigmoid}
  \begin{aligned}
    p_k&=\sigma(z_k)=\frac{1}{1+\exp(-z_k)},\\
    L(x,\bm y)&=-\sum_{k=1}^{K}
    \left[y_k\log p_k+(1-y_k)\log(1-p_k)\right].
  \end{aligned}
\end{equation}
这里每个$p_k$估计“标签$k$存在”的条件概率，彼此不要求加和为1。逐标签二元交叉熵是一种可用的监督目标，并不证明真实标签之间条件独立；共享编码器仍可利用共同线索。若把这些概率相乘来定义完整标签向量的联合分布，才额外采用条件独立的输出假设。两类接口和通用推导见第~\ref{chap:classification}章。

通知中同时出现安排与报名内容，适合输出$(1,1)$，而不是从两类中择一。文体和用途可以另设输出头。有害内容判定也需要这种标签纪律：引用攻击性措辞进行批评、报道威胁和本人发出威胁，不是同一个行为。类别规范应包含对象、作者立场、引用语境及允许证据，关键词最多用于候选筛选。

\subsubsection{多项式朴素贝叶斯文本分类}
\label{sec:nlp-classification-naive-bayes}

当监督样本不多、希望检查每个词怎样影响判断时，可以从类内词频建立基线。固定训练词表$V$，令$x_v$为待分类文档中词$v$的出现次数，$c_{k,v}$为训练集中类别$k$的总词计数，$n_k$为该类文档数。多项式朴素贝叶斯用加性平滑估计
\[
  \widehat\pi_k=\frac{n_k}{n},\qquad
  \widehat\phi_{k,v}
  =\frac{c_{k,v}+\lambda}{\sum_{u\in V}c_{k,u}+\lambda|V|},
  \quad \lambda>0.
\]
比较类别时计算$s_k=\log\widehat\pi_k+\sum_vx_v\log\widehat\phi_{k,v}$，选择最大得分；对所有类别相同的多项式计数系数可以约去。训练保留的是先验和类内词分布，不需要保存整份训练文档。条件独立假设及其与多项式事件模型的关系见第~\ref{sec:classification-naive-bayes}节。

考虑只区分“活动通知”和“其他材料”的教学词表$\{\text{读书会},\text{报名},\text{比赛}\}$。两类的训练文档数为3、2，词计数分别为$(4,3,1)$、$(1,1,6)$。取$\lambda=1$，通知类词概率为$(5,4,2)/11$，其他类为$(2,2,7)/11$。将通知简化成计数$(1,2,0)$，得到
\[
  s_{\text{通知}}=\log0.6+\log(5/11)+2\log(4/11)
  \approx-3.3225,
\]
\[
  s_{\text{其他}}=\log0.4+3\log(2/11)\approx-6.0305.
\]
通知类相对得分高出约2.708，对这两个候选作归一化后约为0.9375。这个数是所选事件模型的后验，不是“原文事实正确”的概率。

平滑使训练中未在某类出现、但已经在词表中的词仍有正概率；词表之外的词则必须事先约定忽略或映射到未知词项，不能临时扩充分母。相同词重复两次会贡献两次对数概率，长文档也会累积更多证据，但模板重复可能造成过强置信。词袋把“无需报名”和“需要报名”中的词当作局部计数，不能自行恢复否定范围。

\subsubsection{线性文本分类：逻辑回归与SVM}
\label{sec:nlp-classification-linear}

词频还可以进入判别式模型，直接学习哪些词区分类别。对训练语料中的文档频率$\operatorname{df}_v$，本章采用$\operatorname{idf}_v=\log((n+1)/(\operatorname{df}_v+1))+1$，将词频乘以该权重后作$L_2$归一化。词表、低频过滤和IDF只在训练语料拟合；验证文档不能参与这些统计，否则输入表示已经使用了留出信息。

例如保留“读书会、报名、比赛”三个维度，通知向量为$\bm x=(0.6,0.8,0)^\top$。某个二分类逻辑回归头取$\bm w=(1,1,-1)^\top$、$b=-0.5$，得分为0.9，故通知概率为$\sigma(0.9)\approx0.7109$。若参考标签为1，损失为$-\log0.7109\approx0.3412$。另一个线性SVM取$\bm w=(0.5,1,-1)^\top$、$b=-0.2$，得分也为0.9；对$y=+1$，铰链损失$\max(0,1-ys)$为0.1。两个0.9都可决定边界一侧，但SVM原始间隔不是概率，不能直接解释为90\%置信。

稀疏表示使得分只需累加非零特征。$L_1$正则化可令部分权重为零，词表过滤则在拟合前删除部分维度；二者发生在不同环节。正类稀少时，可对损失加类别权重或重采样，但部署时仍须在符合目标分布的验证集确定阈值；改变损失权重以后，原始sigmoid也未必保持概率校准。SVM若要输出概率，应额外拟合校准映射，不能把距离线性拉伸后称为后验。通用优化和校准分别回用前卷，此处验收仍以实际标签、实际阈值下的错误为准。

\subsubsection{fastText：词与n-gram表示聚合}
\label{sec:nlp-classification-fasttext}

独立词的线性权重难以在相似短语之间共享统计信息。fastText分类器为词和词$n$-gram建立可学习向量，将特征多重集$F(x)$中的向量平均：
\[
  \bm h(x)=\frac{1}{|F(x)|}\sum_{f\in F(x)}\bm e_f,\qquad
  \bm z=\bm W\bm h+\bm b.
\]
每个训练样本包含文档及类别，交叉熵的梯度同时更新特征表和分类矩阵；运行时仍按同一特征提取规则平均并选类别。原始方法还用哈希桶存储$n$-gram，用层次softmax处理大类别集合\scite{nlp-joulin2017fasttext}

将两句话明确切为“先／报名／再／参加”和“先／参加／再／报名”。本例保留全部四个词特征和三个相邻二元短语，不加入句界特征。第一句的短语为“先报名、报名再、再参加”，第二句为“先参加、参加再、再报名”，因此两份特征多重集的大小都为7。设共同的四个词向量之和为$(2,1)$，第一句三个短语向量之和为$(2,0)$，第二句为$(0,2)$，则均值表示分别为$(4/7,1/7)$和$(2/7,3/7)$。

若“报名先行”标签头为$\bm w=(2,-1)$、$b=0$，两句得分分别为1和$1/7$，sigmoid约为0.7311和0.5357。在教学阈值0.7下才区分出两种安排。例中的特征向量是为说明信息路径而给定的，不是从两句话学出的真实参数；分词、是否加入句界及$n$-gram范围改变以后，分母与参数都应按实际特征集重算。

增加词$n$-gram保留了部分局部顺序，却仍会平均掉长距离结构；哈希碰撞也可能混合无关短语。这里的监督文档分类与fastText字符子词词向量训练不是同一个训练任务，不能把后者处理未登录词的机制直接当成本节分类器已经采用的配置。

\subsubsection{TextCNN：局部模式与池化}
\label{sec:nlp-classification-textcnn}

固定短语表无法枚举所有近义表达。TextCNN让一个滤波器从连续词向量窗口中学习局部分类线索。若词向量$\bm e_j\in\mathbb R^d$，窗口宽为$r$，则一个滤波器产生
\[
  a_j=f\!\left(\bm w^\top
  [\bm e_j;\ldots;\bm e_{j+r-1}]+b\right),\qquad
  h=\max_j a_j.
\]
分号表示拼接。对不同窗口宽和多个滤波器分别最大池化，再将所得特征拼接送入分类头。监督是文档标签；在一次前向计算中，最大池化的梯度主要回到被选中的窗口。原始句子分类方法及多通道变体见Kim的论文，卷积通用原理见第~\ref{chap:convolutional-networks}章\scite{nlp-kim2014textcnn}

将“来宾须提交报名信息”切成五个教学词元，并暂用一维词向量$(0,0.5,1,1,0)$。宽2滤波器取$\bm w=(1,1)$、$b=-0.5$、$f=\operatorname{ReLU}$，四个响应为$(0,1,1.5,0.5)$，最大响应1.5来自“提交／报名”。若报名头为$z=2h-1$，则$z=2$，输出概率0.8808。这个计算把某一窗口与标签连接起来，但它不是一份完整解释证明：其他滤波器、上下文和后续权重也可能改变判断。

把原文改为“无需提交报名信息”，触发窗口仍然存在；否定词若落在窗口之外，最大池化可能保留错误线索。扩大窗口或组合多个窗口能增加可用信息，却不保证理解跨句条件。“校内无需报名；校外须提交信息”还要求按对象组织两条规则，不能只选最强局部响应。长文档分段池化时尤其要检查条件句是否被截断。

\subsubsection{BERT类模型：上下文编码与分类头}
\label{sec:nlp-classification-bert}

否定、对象和条件相互作用时，需要让词元表示随整段语境变化。BERT类模型将词元、位置及输入段信息送入上下文编码器，从分类位置或掩去补齐位置后的池化结果取得$\bm h$，再接式~\eqref{eq:nlp-classification-softmax}或式~\eqref{eq:nlp-classification-sigmoid}。预训练参数提供初始化，训练集的任务标签通过反向传播调整编码器和分类头；不是在预训练词向量之后凭空得到栏目名称。架构与预训练目标见第~\ref{chap:transformer}章。

下面把多标签训练和阈值选择分开算完。只考察“时间安排、报名说明”两个标签，通知参考值为$\bm y=(1,1)$。取$\bm h=(1,2)^\top$、$\bm b=\bm0$，分类矩阵两行为$(\log4,0)$与$(\log1.5,0)$，则$\bm p=(0.8,0.6)$。使用不除以标签数的单样本损失，有
\[
  L=-\log0.8-\log0.6=0.7340,\qquad
  \frac{\partial L}{\partial\bm z}=(-0.2,-0.4)^\top.
\]
为单独观察分类头，暂冻结编码器和偏置。学习率0.1的一步梯度更新是
\[
  \bm W^{(1)}=\bm W^{(0)}
  -0.1(\bm p-\bm y)\bm h^\top.
\]
两行分别增加$(0.02,0.04)$和$(0.04,0.08)$，新得分分别增加0.1和0.2，概率约变为$(0.8155,0.6469)$，损失下降。这一步使用的参考标签只来自训练样本；它既没有选择阈值，也没有证明泛化。

现另给出一个已经训练完成、参数冻结的模型在独立验证集上的得分，见表~\ref{tab:nlp-classification-thresholds}。参考值表示正文是否包含相应说明，而不是活动有没有时间、来宾实际上是否报名。阈值候选网格为$\{0.4,0.5,0.6\}$，约定$p_k\geq t_k$就输出标签。对报名标签，阈值0.4得到两次真阳性、一次假阳性、零次假阴性，故F1为$4/5=0.8$；阈值0.5的F1为0.5，阈值0.6为$2/3$。它的最佳网格阈值因而是0.4。

\begin{table}[htbp]
  \centering
  \begin{tabular}{@{}ccccc@{}}
    \toprule
    验证文档 & 时间得分 & 报名得分 & 时间参考 & 报名参考 \\
    \midrule
    A & 0.80 & 0.60 & 1 & 1 \\
    B & 0.70 & 0.55 & 1 & 0 \\
    C & 0.40 & 0.45 & 0 & 1 \\
    D & 0.30 & 0.35 & 0 & 0 \\
    \bottomrule
  \end{tabular}
  \caption{用于选择阈值的完整教学验证集。它与前面的单步训练样本分开。B只有时间安排却获得偏高的报名分数，C的报名说明得分偏低，使两个标签需要不同的阈值。}
  \label{tab:nlp-classification-thresholds}
\end{table}

时间标签在0.5与0.6下都完全正确，按事先约定的“并列时取较小阈值”规则选0.5。逐标签阈值$(0.5,0.4)$得到宏F1为$(1+0.8)/2=0.9$；把两个标签的计数相加，得到$\mathrm{TP}=4,\mathrm{FP}=1,\mathrm{FN}=0$，微F1为$8/9$。宏平均给予两个标签相同权重，微平均给予每个标签判定相同权重。若只许一个全局阈值，0.4、0.5、0.6的宏F1分别为0.8、0.75、$5/6$，因此按宏F1选择0.6。不能拿0.5这个惯用数值替代验证，也不能在测试集重新挑较好阈值。

新文档若得分为$(0.65,0.48)$，逐标签方案输出两项，全局0.6方案只输出时间安排。两套规则都应连同模型冻结后进入测试。小验证集的最优阈值可能不稳定，罕见标签应记录支持数，必要时使用更保守的统一阈值或共享校准。F1分母为零的类别必须预先约定记零还是排除，不能为提高宏平均临时改变规则。

层次标签还要求恢复合法输出。例如细标签“报名说明”已入选，父标签“活动通知”却未入选，可按规范补入祖先，或者在训练与解码中显式加入约束；补入后的完整标签集合需要重新评价。长文档则先保留分段位置，对段表示平均、最大池化或学习注意权重；平均可能稀释短报名条款，最大池化可能忽略它的适用对象，重叠分段还应避免重复计权。

未见标签不能由固定$K$维分类头直接产生。另一条接口将文档与“包含报名条件”等标签描述配对评分，或受约束地生成合法标签名称；此时需记录描述模板、候选集合和训练时是否见过该标签。候选中得分最高只说明相对择优，不能自动识别“所有候选都不适用”。未知检测应有独立的拒绝条件和未知样本验收；生成标签还要核对解析成功率，不能把非法字符串默默映射到最相近类别。

\subsection{情感与情绪}
\label{sec:nlp-classification-sentiment}

内容分类说明文本在谈什么，评论分析还要说明作者怎样评价所谈对象。\term{情感极性}（sentiment polarity）表示对指定对象的正面、负面或中性评价；\term{情绪}（emotion）表示文本表达的愤怒、喜悦、担忧等状态。情绪可以并存，也可以带强度，两者都不能仅由正负极性唯一恢复。“服务不错，但价格太高”包含两个方向相反的评价；给整篇一个“中性”标签可能只是相互抵消，不能替代方面级分析。

标注可以覆盖整条评论、指定目标或局部方面。参考中还应说明意见持有者，区分作者意见与引述意见；需要定位方面和证据时，沿用前章的跨度与对象绑定。SemEval-2017 Task 4分别设置整条推文情感、目标条件情感以及语料级分布估计，说明输入粒度和输出对象应分别定义，而不是一套极性标签涵盖全部方面分析\scite{nlp-rosenthal2017sentiment}

\subsubsection{情感词典与否定规则}
\label{sec:nlp-classification-sentiment-lexicon}

可检查的词典基线由极性词得分、程度权重和否定范围构成。对目标$a$的意见词集合$O_a$，一种教学规则是
\[
  S(a)=\sum_{j\in O_a}v_jm_jr_j,
\]
其中$v_j$是领域词典极性，$m_j$是程度权重，$r_j$是经句法或局部规则确定的否定修正。这里的加法是设计的评分规则，不是已校准的概率模型。转折前后、引号内外和不同目标之间应先分组，再汇总。

对“价格不算高，但服务很差”，在价格评价语境中取“高”的$v=-1$，“不算”将其修正为$r=-1,m=0.5$，得到价格得分$+0.5$；取“差”的$v=-2$，“很”的$m=1.5$且无否定，得到服务得分$-3$。按零点划分，价格略正面、服务负面。整句求和为$-2.5$，只适用于业务确实要求整体评价的情形。把“不算”改成“并非不高”需要重新判断否定组合，不能永远看最近一个“不”。

规则应保留被命中的词、修正范围和目标，以便核查错误。“服务真好，等了两个小时也没人理”中的表面正词不能直接证明正面；“高”在价格、性能和风险中也不共享固定极性。词典若来自标注数据，扩词与阈值选择只能用训练和开发数据；手写词典则应说明人工依据，不能称其为无假设的方法。

\subsubsection{BERT目标条件情感分类}
\label{sec:nlp-classification-sentiment-bert}

目标条件分类把“评价谁”放进模型输入，例如\code{[CLS]}目标\code{[SEP]}评论\code{[SEP]}。训练样本是$(x,a,y)$，而不只是$(x,y)$；同一评论可以对应多条目标不同的样本。分类头按极性标签作softmax，等级任务可以使用有序目标或等级回归，选择依据来自标注协议。不能将目标位置遗漏后仍声称是在做目标级判断。

对“服务不错，但价格太高”，设标签次序为正面、负面、中性。以“服务”为目标时，教学logits为$(2,0,-1)$，得到$(0.8438,0.1142,0.0420)$；将目标替换为“价格”，logits为$(0,2,-1)$，正负两项交换。两条参考标签分别为正面和负面，相应交叉熵均为约0.1698。这里固定文本、只改变目标，检查的是目标是否真正进入判定；若输出完全不变，应排查模型是否只利用整体语气。

指定目标不一定在原文出现。“这次安排很周到”对地点的评价可能没有证据，不能强制选出正面地点标签。长评论应保留与目标相连的句段和条件，情绪多标签则换用逐类sigmoid并重新标注，强度回归也需要单独的强度参考。目标级准确率、宏F1与证据一致性应分开检查。

\subsubsection{方面与观点的联合抽取模型}
\label{sec:nlp-classification-aspect-opinion}

如果目标事先未知，系统必须先找到方面及观点，再判断它们的关系。令候选方面跨度为$a=(s_a,e_a)$、观点跨度为$o=(s_o,e_o)$，编码器分别计算边界、跨度和配对表示；可用
\[
  s(a,o,c)=g_c([\bm h_a;\bm h_o;\bm h_a\odot\bm h_o])
\]
为极性$c$评分。训练参考包含跨度及配对，目标可由方面损失、观点损失和配对极性损失加权组成；权重在开发集确定。未标配对是否当负例取决于标注是否穷尽，不能把漏标自动当作“明确无关系”。

将“服务／不错／但／价格／太高”视为五个词元，候选方面为位置1、4，观点为位置2、5。设两个正确配对的最高得分为$(1,2,\text{正面}):2.0$、$(4,5,\text{负面}):1.8$，两个交叉配对只有0.1、0.2。在开发集确定的接受门槛1.0下，恢复“服务—不错—正面”“价格—太高—负面”，并把词元位置映射回原文字符跨度。若一个观点同时修饰“场地和服务”，配对允许一对多；若“太贵了”没有显式方面，需要另有隐含方面标签，不能伪造原文跨度。

结构化生成也可输出这些三元组，但生成文本必须回查原文边界和对象绑定。评价时分别报告候选覆盖、方面和观点边界F1、完整三元组F1；已给定方面后的极性正确率，不能代替端到端恢复质量。抽取机制回用第~\ref{chap:nlp-information-extraction}章。

\subsection{目标立场}
\label{sec:nlp-classification-stance}

“改期真让人头疼，不过为了让校外来宾赶上，我支持推迟一天。”语气消极，却对“推迟一天”明确赞成。\term{立场检测}（stance detection）判断文本对指定对象或主张的支持、反对或其他约定关系，输入必须包括目标。情感回答“评价感觉怎样”，立场回答“对这个主张站在哪一边”，两者的证据不同。

目标可能是“活动改期”这样的名词短语，也可能是“应推迟一天”这样的完整主张。对隐含立场，允许使用哪些背景资料必须预先规定；只提供原文时，没有相关讨论可以标为无关，不能由模型自行补足作者态度。Bi-STANCE将中英文、名词短语与主张目标、未见目标及跨域条件区分开，为分开验收这些条件提供了任务依据\scite{nlp-zhao2025bistance}

\subsubsection{BERT目标条件立场分类}
\label{sec:nlp-classification-stance-bert}

将目标和文本组成有序句对，分类头输出支持、反对、中立、无关中的任务定义类别。部分数据集把中立与无关合并，此时输出只说明落入合并类，不能进一步解释成“作者经过权衡保持中立”。训练样本按目标划分时，还应保证同一目标的改写不会同时进入训练与未见目标测试。

对上面的评论，以“应推迟一天”为目标，给定教学logits$(2,0,-1,-2)$，支持概率约0.8310；以“应维持原定日期”为目标，若对应logits为$(0,2,-1,-2)$，则反对概率约0.8310。两次softmax都在各自四个标签内归一化，不在不同目标之间归一化。若文本改为“有人支持推迟，但我没有表态”，引述者和作者的观点应分开；只抓到“支持”会得到错误监督。

已见目标测试检查同类目标上的泛化，目标留出检查新目标，语言留出检查跨语言迁移；这三个条件可能叠加，不能只给一个总体F1。目标名称可揭示训练集偏好，因此还需用目标替换、低词汇重合和引用观点的成组例子检查模型。

\subsubsection{自然语言推断式标签描述匹配}
\label{sec:nlp-classification-stance-nli}

类别数量少而目标常变化时，可以把标签改写成假设：“作者支持将读书会推迟一天”“作者反对将读书会推迟一天”。原文作为前提，假设作为待检验陈述，用自然语言推断模型取得蕴含分数；推断关系的完整定义见第~\ref{sec:nlp-classification-nli}节。模板在此把固定标签变成可输入的文字，并没有免除监督或校准需求。

假设支持、反对、中立三个描述的蕴含概率分别为$(0.75,0.08,0.12)$。这些值来自三个不同的前提—假设三分类，各自有自己的分母，所以不必相加为1。可在开发集确定接受门槛0.6和领先差0.2，先检查最大值及其与次大值之差，再映射回“支持”；此例差为0.63，满足规则。若得分变成$(0.31,0.29,0.28)$，虽仍有最大项，却应保留未定，而不是把相对第一名强行当作支持。

直接立场分类使用目标条件的立场标注；描述匹配可以复用已有NLI模型，但模板的语义、目标否定和领域变化仍需用立场开发集检验。若加入背景资料，训练和测试的可见条件必须一致。“作者感到失望”不是“作者反对改期”的同义标签描述，不能以情感蕴含替换立场证据。

\subsection{文本来源与生成参与判定}
\label{sec:nlp-classification-origin}

一份流畅的通知可能由人工撰写、模型生成，也可能先由人工写草稿再让模型润色。仅问“最后谁动了笔”，会把不同创作过程压成同一个来源标签。\term{生成参与判定}（generation involvement detection）估计文本形成过程中模型参与的类别；作者归属则在规定的候选作者之间判断来源，两者不是同一任务。

Creator-Editor研究（\citeyear{nlp-li2026creator}，见本章文献说明）明确区分纯人工、纯模型、模型润色人工文本和经人工化编辑的模型文本。本章借这一划分说明创作者与编辑者标签不同，不将其结构特征解释为不可伪装的个人身份特征。
实际数据还可能有多轮交替编辑、多人协作或仅局部改写，应按可获得的过程记录标注，必要时增加“混合但过程未知”，而不是为了四分类方便篡改来源。

\subsubsection{文体特征与来源监督分类}
\label{sec:nlp-classification-origin-style}

可复核基线从字符$n$-gram、功能词比例、标点频率和句长统计形成向量，经过标准化后训练线性分类器。作者归属样本需要已知作者、采样时间及题材；生成参与样本还要记录生成系统、提示、编辑方式与程度。目标标签来自来源记录，不能来自另一个检测器的猜测后又当作独立真值。

例如取三个特征为感叹号比例、功能词比例及归一化句长，某通知的向量为$(0.02,0.30,0.50)$。在仅有甲、乙两个候选作者的教学分类器中，甲的logit为$2x_1-x_2+x_3=0.24$，乙为$-x_1+x_2=0.28$，softmax给出约$(0.4900,0.5100)$。若未知来源接受阈值为0.8，这份文本应输出“未能归属”，而不是把乙报告成确定作者。阈值应在留出的未知作者样本上验证；低最大概率规则本身并不保证发现所有未知作者。

改为四种生成参与标签时，需要重新训练四分类头，而不能沿用甲乙得分解释创作过程。上下文编码器也可替换文体表示，但更可能利用题材、平台模板或生成系统专属格式。按主题、作者或生成系统留出能分别暴露这些捷径；混合编辑应按编辑程度分层报告错误，短文本和模板化通知尤其缺少稳定来源信号。统计分类标签不能作为作者身份或违规行为的确定证明。

\subsubsection{困惑度与词元概率检测}
\label{sec:nlp-classification-origin-perplexity}

若能访问某个语言模型的词元概率，可以检查文本是否符合该模型的分布。给定词元序列$w_{1:T}$，平均对数概率和困惑度为
\[
  \bar\ell(x)=\frac1T\sum_{t=1}^{T}\log p(w_t\mid w_{<t}),
  \qquad \operatorname{PPL}(x)=\exp(-\bar\ell(x)).
\]
统计前必须固定模型、词元化、上下文截断和起止符是否计分。不同词元表下的每词元困惑度不能直接比较。若某模型给三个词元的概率为$(0.5,0.25,0.125)$，则平均负对数概率为$\log4$，困惑度为4。

可以用已知来源的开发集选择$\bar\ell$阈值，也可以把概率均值、方差或词元排名组成监督特征。低困惑度只表示该模型容易预测这段文字，常见通知模板、公式化人工文本都可能满足这一点；领域缩写和罕见地点则可能使真实模型输出的困惑度上升。把“海棠楼201”替换成罕见楼名，会改变概率，却不改变来源过程。跨语域、跨生成系统、重写和长度变化都须重新检验，而不能沿用一个固定阈值作身份判断。

\subsubsection{DetectGPT：扰动与概率曲率信号}
\label{sec:nlp-classification-detectgpt}

单看原文概率混合了题材和常用程度的影响。DetectGPT比较原文与附近改写在某个待检验来源模型下的概率差：由扰动模型$q(\widetilde x\mid x)$遮住若干跨度并补写，得到$m$份扰动文本，再用来源模型$p$评分。原始算法将原文相对扰动均值的落差用扰动分数的样本标准差归一化\scite{nlp-mitchell2023detectgpt}

令$\ell(x)$为所采用计分实现的对数概率，$\bar\ell=m^{-1}\sum_j\ell(\widetilde x_j)$，则检测量为
\begin{equation}
  \label{eq:nlp-classification-detectgpt}
  d(x)=\frac{\ell(x)-\bar\ell}
  {\sqrt{\frac{1}{m-1}\sum_{j=1}^{m}
  [\ell(\widetilde x_j)-\bar\ell]^2}}.
\end{equation}
分母为零时不能直接相除，应拒绝该次估计或使用事先约定的稳定化规则。原文以$\log p(x)$说明概率曲率；实际计分若采用平均每词元对数概率，就须对原文和扰动统一使用这一口径，并重新验证阈值，不能混用总概率和平均概率。

为隔离长度影响，取四份等长教学改写，原文总对数概率为$-10$，扰动文本为$(-12,-11,-13,-12)$。均值为$-12$，样本方差为$2/3$，故$d=2/\sqrt{2/3}\approx2.4495$。在开发集固定阈值2.0后，它会被判为偏向该模型生成的样本；若只有三份扰动、或换一个扰动器，估计及阈值都可能改变。改写可以围绕“读书会现改至周五”调整措辞，但不能假定任意扰动都严格保持语义。

该分数依赖来源模型概率可访问、扰动质量及目标分布的假设。它是在离散文本邻域上取得的统计信号，不是直接计算了文本的确定几何曲率，更不是检验出作者身份。对未知来源模型、经过大幅重写或人机共同编辑的文本，应重新定义目标类别，并在对应样本上评价误报和漏报。

\section{文本匹配与关系判定}
\label{sec:nlp-classification-matching}

单篇分类固定了标签表，句对任务还要固定关系及其方向。“机器学习读书会改到周五”与完整通知很相似，但省略了日期、地点和报名条件；它可以是简要概括，却不是保留全部条件的等义复述。完整通知能够支持这句概括，反过来不能恢复全部通知。“周五几点开始？”与“14:00开始”用词不同，后者却比一段详细的读书会介绍更适合回答查询。下面沿这些差别选择监督和计算方式。

\subsection{语义相似度}
\label{sec:nlp-classification-similarity}

\term{语义文本相似度}（semantic textual similarity，STS）给句子对一个连续分数，表示两句话在约定标注规范下有多接近。人工评分可以划分为无关、部分相关、主要内容相同和等义等等级，再求平均；模型既可回归参考分数，也可学习句向量后计算距离。跨语言评分还须记录语言对及标注者能否直接理解两端，不能把机器译文的误差当成原句差异。

独立编码便于预计算候选，联合编码允许两句逐词交互；受约束生成一个数字也是一种接口，但应固定评分区间、非法输出处理和提示模板，并用同一参考集检验。无论采用哪个接口，修改“4月18日”为“4月19日”、将“须报名”改成“无需报名”，都应成为单独的反例组。平均相关性高不保证关键条件变化时会充分降低分数。

\subsubsection{TF-IDF与词向量相似度基线}
\label{sec:nlp-classification-similarity-baseline}

沿用训练语料拟合的词表和IDF，将两句表示为$\bm x,\bm y$，非零向量的余弦为
\[
  s(\bm x,\bm y)=\frac{\bm x^\top\bm y}{\|\bm x\|_2\|\bm y\|_2}.
\]
若只保留“读书会、周五、周六”三个等IDF维度，两句的向量分别是$(1,1,0)$与$(1,0,1)$，余弦为$1/2$。日期词被词表过滤以后，两句都只剩“读书会”，分数反而变成1。空向量应返回不可比较或预先约定的缺省分数，而不能除以零。

词向量平均将词元向量累加后除以有效词数，可使用固定或任务训练的嵌入，也可按IDF加权；计算余弦前仍要处理零范数。“甲通知乙”与“乙通知甲”平均后完全相同，因为求和不保留角色。把句中数字全部映射成一个占位符同样会让人数不同的句子重合。这些基线能够快速找同主题候选，不能单靠高分证明事实相同。

评价时，Pearson相关系数比较预测与参考分数的线性关联，Spearman比较排序，并对并列秩按约定处理。参考为$(1,2,3)$、预测为$(11,12,13)$时，两种系数都为1，但量纲与阈值已经错开。部署若要求$s\geq t$才接受某种关系，还须在开发集选择$t$，在测试集报告该阈值下的误报与漏报。

\subsubsection{Sentence-BERT：孪生编码与句向量}
\label{sec:nlp-classification-sbert}

为了让近义表达在向量空间中接近，Sentence-BERT（SBERT）把同一个上下文编码器用于两句，对有效词元作池化，得到$\bm u=f_\theta(a)$与$\bm v=f_\theta(b)$。共享参数不是把两个输入拼接：每句在编码阶段只能看到自身。原始工作分别讨论句对分类、余弦回归和三元组训练；所选目标必须随模型配置记录\scite{nlp-reimers2019sbert}

例如，均值池化只平均非补齐位置，归一化后的通知向量为$\bm u=(1,0)$，近义句为$\bm v=(0.8,0.6)$，余弦为0.8。若将STS参考分数缩放到$[0,1]$后得到0.9，单对均方误差为$(0.8-0.9)^2=0.01$。梯度经过余弦、池化和共享编码器回传，使这个句对的方向更接近参考要求。运行时关闭训练随机性，编码新句，按同一归一化和评分函数比较。

三元组样本$(a,p,n)$则标明哪个候选应更接近锚点。以欧氏距离和间隔$m$为例，损失为
\[
  L_{\mathrm{triplet}}
  =\max\{0,\|f(a)-f(p)\|_2-\|f(a)-f(n)\|_2+m\}.
\]
若负例向量为$(0.6,0.8)$，$m=0.5$，上述正、负距离为$\sqrt{0.4}$和$\sqrt{0.8}$，损失约0.2380；仅有“正例比负例近”仍未满足间隔。原文句对分类还将$[\bm u;\bm v;|\bm u-\bm v|]$送入分类器，与直接用余弦回归并非同一目标。独立向量便于索引，但序列信息能否保留下来，取决于训练和有限维表示；角色交换与条件遗漏仍需专项测试。

\subsubsection{SimCSE：句向量对比学习}
\label{sec:nlp-classification-simcse}

没有人工相似度分数时，可以从同一句子的两个随机编码视图建立正对。无监督SimCSE对句子$x_i$独立应用两次dropout，得到$\bm h_i,\bm h_i^+$；同批其他句子的第二视图作为负候选。取归一化向量、温度$\tau>0$，一个方向的损失为
\[
  L_i=-\log
  \frac{\exp(s(\bm h_i,\bm h_i^+)/\tau)}
       {\sum_{j=1}^{B}\exp(s(\bm h_i,\bm h_j^+)/\tau)}.
\]
分母包括正对。它要求同句的表示对训练扰动稳定，并与批内其他句子区分；不是用dropout产生了新的人工同义句\scite{nlp-gao2021simcse}

若一个锚点面对正例及两个负例的余弦依次为$(0.8,0.6,0)$，$\tau=0.2$，softmax的未归一化权重为$(e^4,e^3,1)$，正例概率约0.7214，损失约0.3266。只看正例余弦0.8还不够，负例0.6在低温度下仍占据较大竞争权重。偶然同义的批内“负例”会受到排斥，应检查数据重复及采样方式。

监督SimCSE使用NLI中的蕴含句作为正例，矛盾句作为难负例，并将批内候选纳入分母；训练信号于是从“同句不同扰动”变成了已标注的意义关系。它仍输出供相似计算的句向量，并没有使余弦成为有方向的蕴含概率。运行时通常关闭dropout，一句只编码一次；用“甲通知乙／乙通知甲”和“必须报名／无需报名”检验表示，不能把训练目标降低本身当作关系任务验收。

\subsection{复述与语义等价}
\label{sec:nlp-classification-paraphrase}

\term{复述}（paraphrase）是不同表达在任务规定范围内传达相同含义。规范可以容许同义替换和表达省略，但必须说清哪些信息必须保留。本章若按“保留活动事实和适用条件”验收，“校外来宾须在4月16日18:00前提交报名信息”与“校外来宾提交报名信息的时间须早于4月16日18:00”可以作为正例；两句保留相同对象、动作和严格早于该时刻的条件。把“校外来宾”改成“外校人员”未必覆盖同一对象，把“提交报名信息”改成“完成报名”也可能加强动作要求，不能未经约定就标为等义。

删去“校外”会扩大对象，删去时刻会改变截止条件，均应构造负例检查。若业务只需要不保留细节的概括，应另定近似复述规范，不能把严格等价测试中的错误事后解释为“意思差不多”。

正负样本不能只按词汇重合采集。PAWS式反例用高重合但角色或顺序不同的句子暴露这一问题，例如“甲从海棠楼到图书馆”与“甲从图书馆到海棠楼”。分类训练应包含近义低重合正例以及否定、数量、角色、条件变化的难负例。若部署关系要求对称，还须保证交换两句不会改变接受结果。

\subsubsection{双编码器复述判定}
\label{sec:nlp-classification-paraphrase-bi}

在已取得句向量后，可用关系特征
\[
  \bm r=[\bm u;\bm v;|\bm u-\bm v|;\bm u\odot\bm v],
  \qquad p=\sigma(\bm w^\top\bm r+b)
\]
训练二分类头。$\odot$表示逐维乘积。这里增加乘积是本节的关系头选择，不宣称就是原始SBERT分类配置。参考标签$y\in\{0,1\}$给出二元交叉熵$-y\log p-(1-y)\log(1-p)$，训练可更新编码器及关系头，也可冻结编码器。若前两段拼接使用任意不同权重，该头并不天然对称；可约束对应权重相等、只用对称特征，或平均两个输入次序的预测。

给一组对称头的完整计算。取$\bm u=(1,0)$，保留全部条件的改写为$\bm v=(0.96,0.28)$，两向量范数均为1。令拼接原向量的权重为零，差向量的权重为$(-2,-2)$，乘积权重为$(3,3)$，偏置为$-0.5$，则
\[
  z=-2(0.04+0.28)+3(0.96)-0.5=1.74,\qquad
  p=\sigma(1.74)\approx0.8507.
\]
对正例，损失约0.1617；其logit梯度$p-1\approx-0.1493$，将推动该正例得分上升。若删去关键条件的候选向量为$(0.8,0.6)$，同一头得到$z=0.3,p=0.5744$。在独立开发集固定阈值0.8后，前者接受，后者拒绝。这里的向量和阈值是教学给定，真实阈值必须按验证错误选择。

候选向量可以离线保存，但关系特征仍须逐对计算；若头过于复杂，也会增加大规模扫描成本。更根本的限制在表示：若角色相反的两句恰好编码成相同向量$(1,0)$，该头给出$\sigma(2.5)=0.9241$，再好的阈值也难以将它与真实等义句分开。应把这种错误反馈给表示训练、难负例采样或后续交叉编码判定，而不是只提高余弦阈值。

\subsubsection{交叉编码器复述判定}
\label{sec:nlp-classification-paraphrase-cross}

交叉编码器把两句作为\code{[CLS]}句A\code{[SEP]}句B\code{[SEP]}联合编码，让“出发地”“目的地”及其词元在注意力中相互比较，再由句对表示预测二元标签。训练输入是完整句对与参考关系，损失仍是二元交叉熵；它不要求先把每句话压成可独立复用的单向量。联合注意力的架构见第~\ref{chap:transformer}章。

为展示分类头怎样使用交互结果，暂把一个已训练的三维句对表示解释为“角色一致线索、条件缺失线索、否定冲突线索”。这只是可手算的教学坐标，不意味着真实BERT的隐藏维度有这种固定名称。设$z=4h_1-3h_2-5h_3-2$。完整改写的$\bm h=(1,0,0)$给出$p=0.8808$；交换出发地和目的地的$\bm h=(0,0,0)$给出0.1192；遗漏报名条件的$(1,1,0)$给出0.2689；否定冲突的$(1,0,1)$给出0.0474。正例的损失$-\log0.8808=0.1269$，梯度经过分类头回到两句交互。

该例说明需要进入判断的信息，不证明模型一定学到了相应线索。验证应逐项替换实体、数量、否定和截止条件，统计固定阈值下的分类结果。交换输入可以改变联合编码表示，因此对称任务也需双向训练或评分对称化。长度限制还可能把第二句末尾的条件截掉，必须报告截断率和保留策略。

若改用生成模型输出“等价／不等价”，应限制合法标签并记录生成失败；解释文字另行核验，不能用一句流畅理由覆盖错误标签。交叉编码无法为所有句对预存一个统一的联合表示，适合对召回候选做细判；效率取舍不能改变等价关系本身的定义。

\subsection{自然语言推断}
\label{sec:nlp-classification-nli}

\term{自然语言推断}（natural language inference，NLI）给定前提$P$和假设$H$，判断$P$是否支持$H$、是否与$H$不相容，或尚不足以决定。三类通常称为蕴含、矛盾和中立。这里“中立”是证据未定，不是情感中性；三类判断依赖数据集约定的通常语义和背景假设，不是任意形式命题的自动证明。

令前提为本章完整通知。在“同一活动的本次确定开始时间”这一条件下，“读书会本次在4月18日14:00开始”被蕴含；“读书会本次仍在4月17日开始”与更正矛盾；“读书会由张老师主讲”是中立，因为通知未提供主讲人。即使最后一句在现实中碰巧为真，也不能据此前提改标蕴含。把第一条假设与前提交换，短句不再推出报名条件，关系转为中立。

不同评测集改变的也是条件：MultiNLI扩展体裁，ANLI通过对抗式数据收集增加困难实例，XNLI考察多语言迁移评价。使用时应分别报告训练语言、评价语言、是否翻译及体裁，不把多语种成绩当成所有语域上的保证。第~\ref{chap:nlp-question-answering-dialogue}章会在这一接口之外增加证据检索、出处和主张核验要求。

\subsubsection{可分解注意力模型}
\label{sec:nlp-classification-decomposable}

可分解注意力模型把句对判断拆成软对齐、局部比较和聚合。对前提词表示$\bm a_i$及假设词表示$\bm b_j$，前馈网络$F$产生
\[
  e_{i,j}=F(\bm a_i)^\top F(\bm b_j),\qquad
  \widetilde{\bm b}_i=\sum_j
  \frac{\exp(e_{i,j})}{\sum_k\exp(e_{i,k})}\bm b_j.
\]
沿列归一化另得每个假设词对前提的软对齐$\widetilde{\bm a}_j$。比较网络$G$分别处理$[\bm a_i;\widetilde{\bm b}_i]$和$[\bm b_j;\widetilde{\bm a}_j]$，对各侧求和，再用$H$分类。参数由三分类交叉熵监督；基础版本不依赖递归编码，可选句内注意力是额外配置\scite{nlp-parikh2016decomposable}

把前提简化为两个概念向量：“读书会”$(1,0)$和“14:00”$(0,1)$，假设只保留“读书会”$(1,0)$。设两条对齐分数为$(\log3,0)^\top$。按行归一化时，每个前提概念只能对齐唯一假设概念，权重都是1；按列归一化时，假设对前提的权重为$(3/4,1/4)$，取得$(0.75,0.25)$。前提时间词与其对齐对象的绝对差为$(1,1)$，假设词与对齐前提的差为$(0.25,0.25)$。这些不匹配量会进入$G$，不会凭空消失。

为演示恢复，暂令$G$保留绝对差的$L_1$和，两侧聚合值为$A=2,B=0.5$，再给教学分类头三类logits为$(A-B,B-A,-1)=(1.5,-1.5,-1)$，softmax约为$(0.8835,0.0440,0.0725)$，输出蕴含。交换两侧则为$(-1.5,1.5,-1)$，输出中立；这里类别次序为蕴含、中立、矛盾。头的数值用于演示方向可以如何进入判定，真实网络必须从标注学习，不能将“大的一侧包含小的一侧”写成通用逻辑规则。

软对齐没有“找不到对应就自动拒绝”的保证：单候选情况下连时间词也会被迫对齐到活动词。基础词袋式比较对角色交换尤其脆弱；保留未匹配特征和引入上下文，才可能让后续分类器利用这类差别。

\subsubsection{ESIM：上下文编码与局部推断组合}
\label{sec:nlp-classification-esim}

ESIM先用BiLSTM编码每个词在本句中的上下文，再计算句间软对齐。对前提一侧，它形成
\[
  \bm m_i=[\bm a_i;\widetilde{\bm b}_i;
  \bm a_i-\widetilde{\bm b}_i;
  \bm a_i\odot\widetilde{\bm b}_i],
\]
另一侧同样处理。前馈投影以后再用BiLSTM组合局部比较，分别作平均池化和最大池化，拼接后输出三类概率。这个第二层序列模型用于把分散的比较信号组成判断；原文的顺序ESIM与树结构扩展应分开。ESIM的原始工作发表于\citeyear{nlp-chen2017esim}年，完整出处见本章文献说明；BiLSTM原理回用第~\ref{chap:recurrent-networks}章。

比较“甲通知乙报名”与“乙通知甲报名”时，表面词相同，但第一次BiLSTM应有机会使“甲作为通知者”与“甲作为接收者”的表示不同。若某位置的教学编码为$\bm a=(0.8,0.2)$，对齐表示为$\widetilde{\bm b}=(0.3,0.7)$，则差为$(0.5,-0.5)$，积为$(0.24,0.14)$，完整八维输入为$(0.8,0.2,0.3,0.7,0.5,-0.5,0.24,0.14)$。这些信号可与谓词、否定词位置的比较一同进入第二层BiLSTM。

假设组合后两位置输出为$(0.2,0.8)$、$(0.6,0.4)$，平均池化为$(0.4,0.6)$，最大池化为$(0.6,0.8)$。它们与另一句的池化结果共同进入分类头，而非把某个局部差直接等同于矛盾。三分类损失监督全链路；包含“不要求甲通知乙”等否定句时，应检查否定范围、角色互换和多处证据组合。上下文编码增加了可用信息，不保证模型免于只看词重合的捷径。

\subsubsection{BERT类模型与生成式关系判定}
\label{sec:nlp-classification-nli-bert}

联合编码的NLI输入必须保留前提、假设的有序身份，分类头以三分类交叉熵训练。若本章蕴含、矛盾、中立三条假设分别得到logits $(3,0,0)$、$(0,3,0)$、$(0,0,3)$，每条首选标签的概率都为$e^3/(e^3+2)\approx0.9094$，单条正确标签损失约0.0949。交换前提和假设需要重新编码，不能直接调换输出概率；改变输入方向是另一次关系判断。

受约束生成可以输出这三个合法标签之一，监督目标转为标签词元的条件似然；若标签含多个词元，应声明序列评分及结束符规则。自由生成解释则是另一项输出，必须检查它是否引用了前提中的实际内容。NLI头选择“蕴含”并不验证来源可靠，也不负责找到支持该主张的文档。

诊断应比较只给假设的模型、词汇重合基线和完整句对模型，防止标注习惯泄漏标签。“所有校外来宾须报名”不推出“所有参加者都须报名”；“有人报名”也不推出“所有人报名”。把量词和适用对象显式纳入反例，比仅增加随机不相关句更能检验推断任务的薄弱处。

\subsection{查询与候选文本的相关性}
\label{sec:nlp-classification-relevance}

检索中的\term{相关性}是候选对当前信息需求的有用程度。查询—文档、问题—候选答案和需求—证据片段可以共享排序器，但正例标准不同。“读书会几点开始？”的答案片段必须给出本次时间；介绍读书会主题的长文只满足主题相关。证据片段还需支持待核验主张，候选答案则要回答问题，二者不能只凭相似度互换。

设三个教学候选为$d_1$“机器学习读书会介绍”、$d_2$“ML-0418更正：4月18日14:00开始”和$d_3$“原计划4月17日举行”。按当前时间查询，$d_2$是精确答案来源，$d_1$只相关于主题，$d_3$可能误导。索引应保留活动编号、来源与版本，否则召回文本后无法确认它在说哪次活动。

\subsubsection{BM25：词汇相关性基线}
\label{sec:nlp-classification-bm25}

BM25依据查询词在文档中的出现、词项区分度及文档长度计分。为明确本节数值口径，使用非负IDF版本
\[
 \operatorname{BM25}(q,d)=
 \sum_{t\in q}
 \log\!\left(1+\frac{N-\operatorname{df}_t+0.5}
                       {\operatorname{df}_t+0.5}\right)
 \frac{f_{t,d}(k_1+1)}
 {f_{t,d}+k_1(1-b+b|d|/\overline L)}.
\]
这里查询词只计一次，$N,\operatorname{df}_t,\overline L$都来自固定索引；$k_1,b$控制词频饱和与长度修正。完整索引与排序理论属于推荐与搜索部分，本节关注这些统计如何对应语言任务。

取$N=4$、某查询词$\operatorname{df}=2$、$k_1=1.2,b=0.75$且$|d|=\overline L$，该词IDF为$\log2$。出现1次贡献0.6931，出现2次贡献$\log2\cdot4.4/3.2=0.9531$，并未简单翻倍。它能奖励“读书会”的匹配，却不能仅靠词频知道哪个日期已被更正。中文分词、活动编号是否拆开和时间规范化都会改变匹配结果，必须作为索引配置记录。

精确词汇基线适合可追踪的召回；同义问题“何时开始”和“几点举行”可能需要改写或其他召回器。即便$d_1$因大量重复“读书会”排名更高，答案相关标注仍应把$d_2$置于其前，不能按算法当前排序反过来定义真值。

\subsubsection{DPR类双编码器：稠密相关性}
\label{sec:nlp-classification-dpr}

DPR分别编码查询和候选片段，以内积$s(q,d)=f_q(q)^\top f_d(d)$评分；原始方法使用两套BERT编码器及分类位置向量。训练由一个查询、相关片段和负片段组成，负例可来自随机、BM25召回及批内其他正例。一个正例的目标是
\[
  L_q=-\log\frac{\exp s(q,d^+)}
  {\exp s(q,d^+)+\sum_{d^-\in\mathcal N_q}\exp s(q,d^-)}.
\]
难负例应与查询有词汇或主题联系却不满足相关标准，错误地把另一份真正支持材料当负例会产生冲突监督\scite{nlp-karpukhin2020dpr}

若查询向量为$(1,1)$，$d_2,d_1,d_3$的向量依次为$(1,1),(1,0),(0,0)$，内积为$(2,1,0)$。正例概率为$e^2/(e^2+e+1)\approx0.6652$，损失0.4076；梯度推动正例超过当前负例。这是未归一化内积配置，不能不加说明换成余弦后沿用阈值。

部署时离线计算片段向量并绑定文档ID、字符区间、版本和出处，查询在线编码后取top-$k$，再恢复原文。对本查询，top-1含$d_2$，Recall@1为1；若它被分段遗漏，任何后续重排都无法凭空找回。多查询平均召回率、最终回答正确率和证据可靠性应分别报告，第~\ref{chap:nlp-question-answering-dialogue}章再连接完整检索问答过程。

\subsubsection{BERT交叉编码器：候选重排}
\label{sec:nlp-classification-rerank}

重排器只处理已召回的候选，把查询和每个片段联合编码。点式训练可预测二元相关标签；成对训练可最小化$\log(1+\exp[-(s^+-s^-)])$，要求正例高于负例。本节采用点式标签说明输入输出，排序时使用logit或其单调sigmoid，不能把这些独立概率误称为候选集合上的归一化分布。

假设初始列表为$(d_1,d_3,d_2)$，联合模型给出相关logits $(0,-1,2)$，对应概率约$(0.5,0.2689,0.8808)$，排序恢复为$(d_2,d_1,d_3)$。以“能够回答当前开始时间”为唯一相关标准，原列表倒数排名为$1/3$，新列表为1；MRR是跨查询的倒数排名平均。若采用分级相关性，nDCG还需明确等级收益和截断$k$。

训练中的时间难负例要保留“原定”“现改至”等关系词，而不只随机换日期。主题相关、包含答案和能够支持答案是三种不同验收口径。若召回集合只有$d_1,d_3$，重排器仍会输出其中第一名，因此还需要无合适候选的拒绝规则；“第一名”本身并不是“证据足够”。

\subsubsection{ColBERT：词元表示与后期交互}
\label{sec:nlp-classification-colbert}

单向量可能压掉某个稀少但关键的词。ColBERT分别编码两端，却保留归一化的词元向量，再在评分阶段交互。对查询词元$\bm q_i$和文档词元$\bm d_j$，原始MaxSim聚合为
\[
  S(q,d)=\sum_{i=1}^{|q|}\max_j \bm q_i^\top\bm d_j.
\]
每个查询词元找到最匹配的文档词元，得分再求和；不是一对一匹配，同一个文档词元可被多次选择。原始实现还包含查询扩展标记及文档标点过滤，实际复现应随配置固定\scite{nlp-khattab2020colbert}

给定两个查询词元和三个候选词元的相似矩阵
\[
  \bm M=
  \begin{pmatrix}
    0.8&0.2&0.1\\
    0.1&0.3&0.9
  \end{pmatrix},
\]
行最大值为0.8、0.9，故MaxSim为1.7。若另一个候选两行最大值为0.8、0.4，得分为1.2。未归一化平均池化的内积等于矩阵全部元素的均值，本例为0.4；MaxSim保留每行局部最佳匹配，二者不仅尺度不同，压缩信息的方式也不同。平均后再归一化还会产生另一种分数，不能沿用0.4。

候选词元可离线编码，交互在查询到来后进行；交叉编码器则让候选表示随查询变化。ColBERT用相关正负候选训练排序，但局部最大匹配仍可能在同一文档的不同段落找到互不相容的线索。诸如“原时间14:00，新时间15:00”的关系需要编码器正确保留，逐词都有高分并不自动构成一致答案。

\subsection{文本对应、重复与版本关系}
\label{sec:nlp-classification-alignment-duplicates}

对应关系先问两组材料中的哪些单元在表达同一部分内容，近重复先问内容重合多少，版本关系还问哪次编辑继承并改变了哪份材料。英文译文可能把通知一句拆成两句；新通知可能只改一个日期。前者需要一对多对齐，后者需要保存高重合下的实质差异，两者都不能被一个“很相似”标签取代。

\subsubsection{Gale--Church式句长对齐}
\label{sec:nlp-classification-galechurch}

句长对齐利用平行文本大体保序、对应片段长度相关的条件。对源语言片段长度$l_s$和目标语言片段长度$l_t$，Gale--Church式标准化差可写为
\[
  \delta=\frac{l_t-c l_s}{\sqrt{l_s\sigma^2}},
\]
其中$c$为平均长度比例，$\sigma^2$刻画长度波动。依据正态尾概率及对齐类型先验取得代价；再允许$(1,1),(1,2),(2,1),(2,2),(1,0),(0,1)$等转移。语言、字符统计单位和删插模型应在平行开发资料上估计，不能把其他语言的参数直接称为中文经验值\scite{nlp-gale1991alignment}

下面给出可复算的小路径。源文两句长度为$(10,20)$，译文三句长度为$(10,10,10)$，可理解为第二句译成了两句。为了只展示路径机制，教学设$c=1,\sigma^2=10$，匹配代价为$-\log[2(1-\Phi(|\delta|))]$，对$(1,1)$不加罚，$(1,2),(2,1),(2,2)$另加1，单边未对应另加3。这些相对罚值是人为给定，不是原论文拟合值；$\Phi$为标准正态分布函数。

令$D(i,j)$为源端前$i$句与目标端前$j$句的最小代价，$D(0,0)=0$。对允许的$(a,b)$，递推为
\[
 D(i,j)=\min_{(a,b)}
 \left\{D(i-a,j-b)+C(s_{i-a+1:i},t_{j-b+1:j})\right\},
\]
不存在的前缀不可选；单边转移使用未对应罚值，不代入$l_s=0$的长度公式。所需累计表为
\[
 \begin{array}{c|rrrr}
 D(i,j)&0&1&2&3\\ \hline
 0&0&3&6&9\\
 1&3&0&2.1479&5.1479\\
 2&6&2.3935&0.7350&1
 \end{array}
\]
其中$(10,20)$的标准化差为1，匹配项1.1479，故$D(1,2)=1+1.1479$；$(20,10)$的差为$-1/\sqrt2$，匹配项约0.7350。$D(2,1)$还须比较两句合并的$(30,10)$，二对一代价约2.3935，优于先跳过第一句再匹配。回溯$D(2,3)=1$得到第二源句对目标第2、3句，再由$D(1,1)=0$得到第一句一对一。句对必须连同原始位置恢复，不能只输出长度列表。

长度相近并不证明翻译正确，重复模板、摘要式翻译和大段重排都可能破坏假设。句对齐用于建立候选对应；词对齐另行输出词之间的联系，不是把这一条动态规划路径细分就能自动得到。

\subsubsection{跨语言表示匹配与全局对齐}
\label{sec:nlp-classification-crosslingual-align}

跨语言句编码器把不同语言映射到可比较空间，以翻译句对或其他跨语监督训练；每句检索top-$k$目标候选后，再加入保序、互斥或位置距离条件。其分数尺度需要按语言对验证，中文同语种阈值不能直接作为中英对应阈值。没有合适对应时应允许未配对。

例如源端两句和目标端两句的余弦矩阵为
$\bigl(\begin{smallmatrix}0.90&0.85\\0.89&0.10\end{smallmatrix}\bigr)$。
各行独立取最大值会都选目标第1句。若只要求一对一而不保序，最大总分配对为交叉的$0.85+0.89=1.74$；若材料确定保序，则不能采用该路径。设每次接受配对的收益为$s-0.8$，跳过单元收益为零，保序动态规划选第一对0.90，第二句未配对，总收益0.10，高于接受低质量第二对。这说明约束和未配对规则都是输出含义的一部分。

一对多情形可把相邻目标句组合后重编码，作为额外候选，再用相同顺序约束恢复路径。应分别测候选覆盖率和最终对应正确率：正确翻译未进候选，是召回问题；已在候选却被全局约束挤掉，是评分或约束问题。IBM模型或fast\_align可提供词对齐证据，但词对应需要聚合才能辅助句匹配，不能把词链接直接当作句级结果。

\subsubsection{MinHash：片段集合的近重复检索}
\label{sec:nlp-classification-minhash}

精确重复可以先规范化允许忽略的空白、编码及标点，再比较哈希，并在需要时核对规范化正文。不能为了增加命中而删去日期或编号。近重复则把连续字符或词片段组成集合$A,B$，用Jaccard相似度$J(A,B)=|A\cap B|/|A\cup B|$描述共享比例；集合不保留一个片段重复出现的次数。

MinHash对全集施加理想随机排列，取集合中排列最靠前的元素作为签名。两个集合最小元素相同，当且仅当并集最靠前的元素落在交集中，因此碰撞概率恰好为Jaccard相似度。多次独立排列的碰撞比例给出估计；工程上常用多组哈希近似排列，性质需结合所用哈希族理解\scite{nlp-broder1997resemblance}

取$A=\{a,b,c\}$，$B=\{a,b,d\}$，对应三种教学通知片段，其中$c,d$含不同日期。四个排列依次为
$(a,b,c,d)$、$(b,d,a,c)$、$(c,a,d,b)$、$(d,c,b,a)$。
两份签名分别是$(a,b,c,c)$与$(a,b,a,d)$，相同两位，估计为$2/4=0.5$，也等于本例精确Jaccard值$2/4$。这是小例恰好一致，不是四个签名总能精确还原；真实检索用签名筛选候选后，应计算原集合相似或核读正文。

当规范明确“某通知主要被另一份包含”时，还可以计算有方向的包含率$|A\cap B|/|A|$，不能与Jaccard混用。短文档片段少，单个日期变化可能大幅改分；长模板共享大量片段，又可能掩盖关键事实变化。因此近重复关系只支持候选合并建议，不能据此直接删除新版本。

\subsubsection{SimHash与版本差异核验}
\label{sec:nlp-classification-simhash}

SimHash保留加权特征对向量方向的影响。理想随机超平面方法对每个随机向量$\bm r_j$记录$\operatorname{sign}(\bm r_j^\top\bm x)$；两个向量夹角为$\theta$时，某位不同的概率为$\theta/\pi$。工程文本指纹常把特征哈希到若干正负位，按词权重累加后取符号。它近似角度关系，与MinHash针对集合交并比例不同，也不能把短指纹的汉明距离当成精确余弦\scite{nlp-charikar2002similarity}

给三个特征的四位正负码：“读书会”为$(+,+,-,-)$，“报名”为$(+,-,+,-)$，“周五”为$(-,+,+,-)$，权重分别为2、1、1。位和为$(2,2,0,-4)$，约定零取正，得到指纹\code{1110}。把日期特征换成码为$(-,-,+,+)$的“周六”，位和变成$(2,0,0,-2)$，指纹仍为\code{1110}，汉明距离为0。例子不是在证明方法无效，而是直接展示短指纹会丢掉实质变化。

候选命中以后要恢复原文差异、实体、数字、否定和条件。共同通知的新旧版本即使高度近重复，也须保存“4月17日”变为“4月18日14:00”的编辑；若没有发布时间、活动编号或编辑链记录，还不能仅凭文本推断谁先谁后。候选覆盖、最终重复判定与版本变更核验应分别评价，才能知道漏的是相近材料还是关键修改。

\section{文本集合的组织}
\label{sec:nlp-classification-organization}

前两节的类别和关系都有外部定义。面对还没有目录的一批通知，问题变成了怎样发现可供浏览的结构。主题模型通常给出词的组织及文档的混合权重；聚类把文档分组；关键词为单篇或一组材料提供导航入口。它们可以协作，却不共享同一个真值：一个簇名概括了多数材料，不意味着每个成员的每句话都符合该名称。

\subsection{主题发现}
\label{sec:nlp-classification-topics}

\term{主题发现}（topic discovery）从一批文本中寻找反复出现的内容模式，并尝试用词语和代表文档解释。这里的主题是模型对共现或表示结构的归纳，不是已经标注好的“报名说明”类别。若事先规定类别并提供标签，应回到监督分类，而不是把预测类别称为发现主题。

令$x_{d,v}$为文档$d$中词$v$的计数，$D$为文档数，$K$为待拟合的主题数，$V$为固定词表。分词、停用词、词形归并、低频过滤和数字处理先定义了模型能看见什么；把“无需”删除以后，词袋方法已无从利用这个否定词。长文档贡献更多词计数，模板段落可能压过短小的实质修改，因此还应检查重复文本和分段单位。训练词表冻结后，新文档的未登录词须按既定规则忽略或映射；空词袋应保留不可推断状态。

下文固定三个教学词“讲座、报名、场地”，训练文档$d_1,d_2$的计数矩阵为
\begin{equation}
 \label{eq:nlp-classification-topic-counts}
 \bm X=
 \begin{pmatrix}2&1&0\\0&1&2\end{pmatrix}.
\end{equation}
$d_1$围绕讲座及报名，$d_2$围绕报名及场地；另留一篇三项各出现一次的混合文档$d_3=(1,1,1)$检查新文档使用。它们是从通知集合压缩出的计数，不是完整自然语言。实际解释应返回原文，确认高权重词是否确实属于同一内容，而不是偶然共现。

表~\ref{tab:nlp-classification-topic-methods}比较四类实现。矩阵分解侧重重构，经典概率模型显式解释词的混合生成，神经主题模型用网络进行参数化或推断，表示聚类则先分文档再找代表词。这个分类用于安排实现路径；神经模型也可以是概率生成模型，两者并不互斥。

\begin{table*}[htbp]
 \centering
 \begin{tabularx}{\linewidth}{@{}p{15mm}p{19mm}Xp{25mm}Xp{23mm}@{}}
  \toprule
  方法 & 输入 & 拟合与输出 & 新文档处理 & 权重的含义 & 评价条件 \\
  \midrule
  LSA & TF-IDF等加权矩阵 & 截断SVD；词方向和文档坐标 & 投影到固定右奇异向量 & 有正有负的坐标，不是概率 & 重构、稳定性、语义用途 \\
  NMF & 非负词权重 & 非负重构；成分词权重与文档系数 & 固定成分，求非负系数 & 归一化系数也不自动成为后验 & 指定损失、重复成分与解释 \\
  pLSA & 词袋计数 & EM；主题词分布及文档比例 & 固定词分布，局部拟合比例 & 在该词生成模型中的混合比例 & 留出词预测及主题解释 \\
  LDA & 词袋及先验 & 后验推断；词主题指派与分布 & 给新文档作后验推断 & 文档主题比例的后验估计 & 同词表似然、稳定性与解释 \\
  NVDM／ProdLDA & 词袋计数 & 推断网络与词重构 & 编码后取得潜在分布 & 前者为高斯坐标；后者为词logits组合权重 & 重构协议、先验及主题重复 \\
  CTM & 词袋与上下文向量 & 联合输入推断；重构词袋 & 同一编码与词表接口 & 取决于解码器，不是类别概率 & 截断、覆盖、解释与稳定性 \\
  BERTopic & 文档向量与词袋 & 降维聚类；簇及代表词 & 固定变换及聚类预测接口 & 成员得分不是词生成混合比例 & 离群率、簇稳定性与代表词 \\
  \bottomrule
 \end{tabularx}
 \caption{主题发现方法的实现与输出。比较基于相同教学材料，未给出性能排名。单篇文档的主题比例、潜在坐标和聚类成员分数只有在各自模型定义内才有意义；不能因数值都处于$[0,1]$就互相替换。}
 \label{tab:nlp-classification-topic-methods}
\end{table*}

\FloatBarrier
\subsubsection{LSA：截断奇异值分解与潜在语义}
\label{sec:nlp-classification-lsa}

\term{潜在语义分析}（latent semantic analysis，LSA）将文档—词加权矩阵近似为低秩结构$\bm X\approx\bm U_r\bm\Sigma_r\bm V_r^\top$。每列$\bm V_r$是一条词空间方向，文档坐标为$\bm X\bm V_r=\bm U_r\bm\Sigma_r$；共同方向把分散词项关联起来。通常先用TF-IDF降低通用词影响，SVD及最优低秩逼近原理见第~\ref{chap:dimensionality-reduction}章。

为了单独看清分解，本例给式~\eqref{eq:nlp-classification-topic-counts}统一词权重1，不作行归一化；这不是按前节公式从这两篇文档估出的IDF。它的两个非零奇异值为$\sqrt6$和2，对应右方向可以取
\[
 \bm v_1=\frac{(1,1,1)^\top}{\sqrt3},\qquad
 \bm v_2=\frac{(1,0,-1)^\top}{\sqrt2}.
\]
只保留第一维时，两篇坐标都是$\sqrt3$，重构矩阵两行都是$(1,1,1)$，只留下“都涉及这组通知词”的共同方向。保留第二维后，坐标分别为$(\sqrt3,\sqrt2)$与$(\sqrt3,-\sqrt2)$，讲座与场地的差别才被恢复。

新文档$\bm x_{\mathrm{new}}$沿固定词表加权后右乘$\bm V_r$，不重新拟合旧文档；$d_3=(1,1,1)$投影为$(\sqrt3,0)$。解释某方向时可以分别列正、负载荷高的词，并核读对应文档。第二方向包含负数，翻转整列符号也不改变分解，因此它不是“两种业务主题的概率差”的天然定义，更不能直接送入需要非负比例的接口。

\subsubsection{NMF：非负成分与主题权重}
\label{sec:nlp-classification-nmf}

若希望每篇文本由非负成分相加解释，可以取$\bm X\approx\bm W\bm H$，其中$\bm W\in\mathbb R_{\geq0}^{D\times K}$为文档系数，$\bm H\in\mathbb R_{\geq0}^{K\times|V|}$为成分词权重。以平方损失$\frac12\|\bm X-\bm W\bm H\|_F^2$为例，从非负初始化交替更新两矩阵，至目标相对变化小于容差或达到迭代上限。乘法更新可从随机正数开始，避免整列零初始化造成无法恢复的成分；分母为零须另定稳定化或重启规则。通用更新和其他损失见第~\ref{sec:dim-nmf}节。

对同一$\bm X$，取
\[
 \bm W=\begin{pmatrix}1&0.5\\0.5&1\end{pmatrix},\qquad
 \bm H=\begin{pmatrix}1&1&0\\0&1&1\end{pmatrix}.
\]
当前重构两行为$(1,1.5,0.5)$、$(0.5,1.5,1)$，损失为1.5。暂固定$\bm H$，按元素乘除更新
$\bm W\leftarrow\bm W\odot(\bm X\bm H^\top)/(\bm W\bm H\bm H^\top)$，
得到
$\bm W'=\bigl(\begin{smallmatrix}1.2&0.25\\0.25&1.2\end{smallmatrix}\bigr)$，
损失降至0.905。接着才轮到$\bm H$，不能把这半轮说成已经收敛。

两成分的高权重词分别为“讲座、报名”和“报名、场地”。新文档固定$\bm H$后求非负最小二乘系数；对$d_3$，最优系数为$(2/3,2/3)$，重构为$(2/3,4/3,2/3)$。把系数再归一化为$(1/2,1/2)$便于展示，但这是人为消除尺度，不是贝叶斯后验。$\bm W$某列乘常数、$\bm H$对应行除同一常数，不改变重构，解释前必须固定尺度。

非负性不保证每个主题独特，也不自动保证稀疏。不同初始化可能得到重复成分或不同分解，稀疏性还可由惩罚显式控制。若改用广义KL损失，计数误差的权重随之改变；不能把平方重构最优解直接解释成pLSA最大似然。

\subsubsection{pLSA：概率潜在语义分析}
\label{sec:nlp-classification-plsa}

pLSA将上述“成分”约束成概率分布。文档$d$有主题比例$\theta_{d,k}=p(z=k\mid d)$，主题$k$有词分布$\phi_{k,v}=p(w=v\mid z=k)$，分别满足$\sum_k\theta_{d,k}=1$和$\sum_v\phi_{k,v}=1$。每个词先按文档比例选择主题，再按主题选择词，所以
\[
 p(v\mid d)=\sum_{k=1}^{K}\theta_{d,k}\phi_{k,v},\qquad
 \ell=\sum_{d,v}x_{d,v}\log\!\left(\sum_k\theta_{d,k}\phi_{k,v}\right).
\]
词袋省略词序，固定计数下与参数无关的多项式系数不影响拟合。原始方法及概率结构见前卷；这里完整走通计数怎样变成文档主题结果\scite{pgm-static-directed-hofmann1999}

对式~\eqref{eq:nlp-classification-topic-counts}取两个主题，初始化为
\[
 \bm\Theta^{(0)}=
 \begin{pmatrix}0.8&0.2\\0.2&0.8\end{pmatrix},\qquad
 \bm\Phi^{(0)}=
 \begin{pmatrix}0.6&0.3&0.1\\0.1&0.3&0.6\end{pmatrix}.
\]
所有元素为正，避免某主题从开始就永远得不到某词。E步计算“当前模型认为这个已观测词来自哪个主题”的责任度
\[
 r_{d,v,k}=\frac{\theta_{d,k}\phi_{k,v}}
                   {\sum_j\theta_{d,j}\phi_{j,v}}.
\]
例如$d_1$中“讲座”来自主题1的未归一化质量为$0.8\times0.6=0.48$，来自主题2为$0.2\times0.1=0.02$，故责任度为$(0.96,0.04)$。它出现两次，贡献的是$(1.92,0.08)$而不是一份责任度。完整非零计数见表~\ref{tab:nlp-classification-plsa-em}。

\begin{table}[htbp]
 \centering
 \begin{tabular}{@{}ccrcccc@{}}
  \toprule
  文档 & 词 & 次数 & $r_1$ & $r_2$ & $x r_1$ & $x r_2$\\
  \midrule
  $d_1$ & 讲座 & 2 & 0.96 & 0.04 & 1.92 & 0.08\\
  $d_1$ & 报名 & 1 & 0.80 & 0.20 & 0.80 & 0.20\\
  $d_2$ & 报名 & 1 & 0.20 & 0.80 & 0.20 & 0.80\\
  $d_2$ & 场地 & 2 & 0.04 & 0.96 & 0.08 & 1.92\\
  \bottomrule
 \end{tabular}
 \caption{pLSA一次完整E步。零计数项不贡献加权计数，无须为其新增观测。每行责任度和为1，加权计数和等于该词的出现次数。}
 \label{tab:nlp-classification-plsa-em}
\end{table}

M步按文档和主题分别归一化：
\[
 \theta_{d,k}^{(1)}=\frac{\sum_vx_{d,v}r_{d,v,k}}{\sum_vx_{d,v}},
 \qquad
 \phi_{k,v}^{(1)}=\frac{\sum_dx_{d,v}r_{d,v,k}}
                      {\sum_{d,u}x_{d,u}r_{d,u,k}}.
\]
$d_1$的主题计数为$(2.72,0.28)$，$d_2$为$(0.28,2.72)$；主题1的词计数为$(1.92,1,0.08)$，主题2为$(0.08,1,1.92)$，各自总数都是3。因此
\[
 \bm\Theta^{(1)}=
 \begin{pmatrix}68/75&7/75\\7/75&68/75\end{pmatrix},\qquad
 \bm\Phi^{(1)}=
 \begin{pmatrix}16/25&1/3&2/75\\2/75&1/3&16/25\end{pmatrix}.
\]
这就是完整一轮EM，而非只更新词分布。将新旧参数代回同一个$\ell$，得到对数似然从$-5.1805$上升到$-4.3572$；重复两步直至相对增量足够小或达到上限，再比较多个初始化。零计数文档没有可估计比例，应跳过或另定缺省规则；若某主题的总责任度为零，其词分布更新分母也为零，需按预定策略重启或去除该主题。

新文档没有训练时的$\theta_d$参数。固定已学$\bm\Phi^{(1)}$，给新$\bm\theta$正初始化，只交替计算该文档责任度和更新该比例，这称为局部拟合。对$d_3=(1,1,1)$从$(0.5,0.5)$开始，主题1三个词责任度为$(0.96,0.5,0.04)$，平均仍为0.5；对只含“讲座、报名”的新计数$(1,1,0)$，第一步则得到$(0.73,0.27)$，还应继续迭代。两者都不能修改训练词分布来迎合新输入。

pLSA为每篇训练文档保存自由参数，没有由此学到“新文档比例应长什么样”的文档先验。另一方面，它与非负分解有关：若令$\widehat x_{d,v}=N_d\sum_k\theta_{d,k}\phi_{k,v}$，在固定文档长度$N_d$及上述归一化约束下，广义KL散度与负计数对数似然只差参数无关项。这个关系不适用于任意缩放的平方损失NMF。

\subsubsection{LDA：潜在Dirichlet分配}
\label{sec:nlp-classification-lda}

pLSA可以为新文档局部拟合，但短文档的自由比例容易受少量词支配。\term{潜在Dirichlet分配}（latent Dirichlet allocation，LDA）给文档比例加Dirichlet先验$\bm\theta_d\sim\operatorname{Dir}(\bm\alpha)$，让稀少观测与跨文档共享的先验共同决定后验。此处LDA是主题模型，不是线性判别分析；原始LDA将主题词分布作为全局参数，下面采用再给主题词分布加Dirichlet先验的常用贝叶斯形式\scite{blei2003}

具体地，取$\bm\phi_k\sim\operatorname{Dir}(\bm\beta)$，每个词位置先抽$z_{d,n}\sim\operatorname{Cat}(\bm\theta_d)$，再抽$w_{d,n}\sim\operatorname{Cat}(\bm\phi_{z_{d,n}})$。得到的主题比例控制词主题指派，不是业务类别出现的概率。概率结构见第~\ref{sec:pgm-static-directed-hierarchy-membership}节；本节选折叠Gibbs采样展示任务中的推断。

令$n_{d,k}$为文档$d$分给主题$k$的词数，$n_{k,v}$为主题$k$中词$v$的计数，$n_k=\sum_vn_{k,v}$为主题总词数。对称先验$\alpha,\beta$下，把当前位置$i$的旧主题指派从所有计数中移出，再计算
\begin{equation}
 \label{eq:nlp-classification-lda-gibbs}
 p(z_i=k\mid\bm z_{-i},\bm w)\ \propto\
 (n_{d_i,k}^{-i}+\alpha)
 \frac{n_{k,w_i}^{-i}+\beta}{n_k^{-i}+|V|\beta}.
\end{equation}
前一项偏好当前文档已使用的主题，后一项偏好能够生成当前词的主题；$|V|$是词表大小，上标$-i$表示移出当前位置，省略的文档长度分母对所有$k$相同。完整采样原理回用第~\ref{chap:pgm-sampling}章。

继续使用两篇文档和三个词。将$d_1$的两次“讲座”初始指派给主题1、一次“报名”给主题2；将$d_2$三词都给主题2。此时主题词计数为$(2,0,0)$和$(0,2,2)$。取$\alpha=\beta=0.5$，准备重采样$d_1$的“报名”：先移出它，文档计数变成$(2,0)$，全局主题词计数变成$(2,0,0)$和$(0,1,2)$。两主题未归一化权重为
\[
 u_1=(2+0.5)\frac{0+0.5}{2+1.5}=\frac5{14},\qquad
 u_2=(0+0.5)\frac{1+0.5}{3+1.5}=\frac16.
\]
归一化得$(15/22,7/22)$。若本次均匀随机数为0.6，落在第一主题区间，就把该词加回主题1；文档计数变为$(3,0)$，全局主题词计数为$(2,1,0)$和$(0,1,2)$。若不先移出旧计数，当前词会用自己给自己增加证据。

按这个状态计算的条件后验均值为
\[
 \begin{aligned}
   \widehat{\bm\theta}_{d_1}
   &=\frac{(3+0.5,0+0.5)}{3+1}=(0.875,0.125),\\
   \widehat{\bm\phi}_1&=(5/9,1/3,1/9),\\
   \widehat{\bm\phi}_2&=(1/9,1/3,5/9).
 \end{aligned}
\]
它只是一个状态下的估计，不是整条链已经收敛。实际要遍历全部词、多轮采样、丢弃预热阶段，并对保留状态的分布估计求平均；主题编号的置换对词生成概率无影响，跨运行比较前须对齐主题。

新文档可以在固定主题词分布下只采样局部指派。以刚才的$\widehat{\bm\phi}$作为教学固定值，给$d_3$的“讲座、报名、场地”初始指派$(1,2,2)$。移出“讲座”后，新文档计数为$(0,2)$，条件权重为
\[
 u_1=(0+0.5)\frac59=\frac5{18},\qquad
 u_2=(2+0.5)\frac19=\frac5{18}.
\]
两者概率各0.5；随机数0.4使该词回到主题1，此状态的文档比例估计为$(1.5,2.5)/4=(0.375,0.625)$。继续更新另外两个词并多轮采样以后才汇总；若四个保留状态的主题计数依次为$(1,2),(2,1),(1,2),(2,1)$，则平均比例为$(0.5,0.5)$。后一个数只是给定样本的平均示例，不能拿一次更新冒充完整后验。

固定$\widehat{\bm\phi}$是新文档的近似使用方式，也可对训练后验中的多组$\bm\phi$重复推断以传播不确定性。变分推断则用可计算分布近似后验，是另一条求解路线。较小的$\alpha$通常偏好更稀疏的文档比例，较小的$\beta$偏好更集中的主题词分布；具体结果还依赖主题数和数据，不能由先验大小直接断言某文档只含一个主题。

受监督主题模型增加文档标签，并使标签似然约束潜在主题及预测；动态主题模型增加时间信息，估计随时间变化且相互关联的主题词分布。二者改变了输入和估计对象，不是给静态LDA输出重新命名。对于极短通知、领域缩写和多义词，先验能缓和稀疏性，却不能恢复被词袋忽略的顺序和条件。

\subsubsection{NVDM：神经变分文档模型}
\label{sec:nlp-classification-nvdm}

逐篇迭代推断之外，还可以训练一个网络，直接从词袋近似其潜在分布。神经变分文档模型（NVDM）由计数$\bm x$预测$\bm\mu(\bm x)$与$\log\bm\sigma^2(\bm x)$，构造
$q(\bm h\mid\bm x)=\mathcal N(\bm\mu,\operatorname{diag}\bm\sigma^2)$。
用$\bm h=\bm\mu+\bm\sigma\odot\bm\epsilon$、$\bm\epsilon\sim\mathcal N(\bm0,\bm I)$重参数化，再以线性logits的softmax重构词分布$\bm p(\bm h)$。训练最大化
\[
 \mathcal L=
 \mathbb E_q\!\left[\sum_vx_v\log p_v(\bm h)\right]
 -D_{\mathrm{KL}}(q(\bm h\mid\bm x)\,\|\,\mathcal N(\bm0,\bm I)).
\]
第一项奖励解释观测词，第二项约束潜在分布；VAE推导见第~\ref{sec:pgm-generative-vae}节\scite{nlp-miao2016nvdm}

对$d_1=(2,1,0)$，给定推断网络输出$\bm\mu=(0.5,0)$、$\bm\sigma^2=(1,1)$。取一份教学噪声$\bm\epsilon=\bm0$展示计算，解码logits设为$(h_1,0,-h_1)=(0.5,0,-0.5)$，词概率约为$(0.5065,0.3072,0.1863)$。这一次重构负对数似然为$-2\log0.5065-\log0.3072\approx2.5408$，高斯KL为$\frac12(0.5^2)=0.125$，负ELBO的单样本估计为2.6658。零噪声只用于算例，训练不能因此取消随机采样的期望。

对解码logits的重构梯度是$N\bm p-\bm x$，本例约为$(-0.4806,-0.0784,0.5590)$，因而应提高“讲座、报名”相对“场地”的概率。新文档使用同一推断网络，不必重新学习每篇自由参数；可用均值表示或按后验采样预测。$\bm h$可为负、各坐标不加和为1，不能称为LDA主题比例。用某坐标的高解码权重词解释时，还应检查相关文档和不同运行，避免把任意坐标轴当成唯一语义主题。

\subsubsection{ProdLDA：摊销推断与专家乘积}
\label{sec:nlp-classification-prodlda}

ProdLDA保留“由网络近似推断”的接口，但将潜在logits映射到单纯形上的文档权重。它在softmax参数空间以logistic-normal形式近似Dirichlet先验，推断网络给高斯参数，经重参数化取得$\bm h$，再令$\bm\theta=\operatorname{softmax}(\bm h)$。词袋重构与近似先验的KL共同训练网络和主题参数，不能因为有softmax就忽略先验约束\scite{nlp-srivastava2017prodlda}

关键区别在词概率的组合位置。LDA使用$\bm p_{\mathrm{mix}}=\sum_k\theta_k\bm\phi_k$；ProdLDA先组合主题词logits$\bm b_k$，再作
\[
 p_{\mathrm{prod},v}
 =\frac{\exp(\sum_k\theta_k b_{k,v})}
 {\sum_u\exp(\sum_k\theta_k b_{k,u})}.
\]
若为比较而取$b_{k,v}=\log\phi_{k,v}$，便得到归一化加权几何平均$\prod_k\phi_{k,v}^{\theta_k}$，而非算术平均。这是“专家乘积”在此处的含义，不是逐主题先softmax再相加。

取主题词分布$(0.6,0.3,0.1)$和$(0.1,0.3,0.6)$，文档权重各0.5。LDA混合得到$(0.35,0.30,0.35)$；ProdLDA对应的未归一化质量为$(\sqrt{0.06},0.3,\sqrt{0.06})$，归一化约为$(0.3101,0.3798,0.3101)$。两个主题都给予“报名”中等质量，乘积相对抑制只受单个主题偏好的两端词，因此“报名”排到第一。这个例子展示机制差异，不宣称它在所有数据上更可解释。

新文档先取得潜在分布，再采样或用约定的确定性近似形成权重和词预测。虽然$\bm\theta$非负且和为1，它在这个解码器里是logits组合权重，并非LDA那样先抽离散主题再抽词的同一生成比例；也不是业务标签后验。应核查主题词重复、后验使用程度和随机运行差异，防止多个主题学成相同词分布。

\subsubsection{CTM：上下文文档表示与主题推断}
\label{sec:nlp-classification-ctm}

词袋推断网络无法区分词相同而顺序不同的材料。上下文主题模型（contextualized topic model，CTM）的CombinedTM配置同时使用词袋和SBERT文档向量：先把上下文向量投影到词表维度，与词袋拼接，送入推断网络得到潜在分布参数，再用ProdLDA式权重和词解码重构词袋。上下文输入改善的是推断线索，不是直接输出了一组可解释主题\scite{nlp-bianchi2021ctm}

将“甲通知乙”与“乙通知甲”简化为相同词袋$(1,1,1)$，给定上下文向量分别为$(1,0)$与$(0,1)$。教学投影
$\bm P=\bigl(\begin{smallmatrix}1&0\\0&1\\1&-1\end{smallmatrix}\bigr)$
产生$(1,0,1)$与$(0,1,-1)$，与词袋拼接后分别为$(1,1,1,1,0,1)$和$(1,1,1,0,1,-1)$。若推断均值相应为$(1,0)$和$(0,1)$，在零噪声展示下，主题权重为$(0.7311,0.2689)$和$(0.2689,0.7311)$。这说明相同词袋允许得到不同推断输入，并不证明角色变化一定应该划成不同主题。

训练仍以观测词袋重构及先验约束为目标；如果这些词的生成分布不需要区分角色，主题也可能保持相同。新文档必须使用同一SBERT、投影和词表，长文档截断后遗漏的报名条件不会因拼接而恢复。验收应记录上下文截断、未登录词比例、重复主题及多次运行稳定性，并同时查看高权重词和代表文档。

\subsubsection{BERTopic：表示聚类与c-TF-IDF}
\label{sec:nlp-classification-bertopic}

若主要需求是把文档分组并提供组名，可以先在上下文表示上聚类。BERTopic以SBERT等编码器取得文档向量，经UMAP降维后由HDBSCAN寻找密度簇及离群文档，再把每个簇内文本合并成一个词袋，用c-TF-IDF选择代表词。向量用于决定文档近邻，词袋用于解释已得到的簇，这两个环节可以分别调整。其\citeyear{nlp-grootendorst2022bertopic}年原始报告的完整出处见本章文献说明。

为走通接口，假设拟合后的结果将$d_1$与补充的讲座通知$(2,0,0)$归入簇A，将$d_2$归入簇B，混合文档$d_3$暂标离群。这组指派是教学给定的聚类输出，不声称由前面三维计数直接运行UMAP得到。暂不让离群项参与代表词统计，簇计数为$A=(4,1,0)$、$B=(0,1,2)$，平均簇长度$\overline L_c=4$。采用簇内$L_1$归一化词频时，
\[
 w_{c,v}=\frac{x_{c,v}}{\sum_u x_{c,u}}
       \log\!\left(1+\frac{\overline L_c}{\sum_{c'}x_{c',v}}\right).
\]
分母的词频是跨簇总次数，不是“出现该词的簇数”，不能直接替换成普通IDF。

三词的逆频率项为$(\log2,\log3,\log3)$，相应的代表词权重为
\[
  \begin{aligned}
    \bm w_A&\approx(0.5545,0.2197,0),\\
    \bm w_B&\approx(0,0.3662,0.7324).
  \end{aligned}
\]
从“讲座、报名”和“场地、报名”可以提出簇名称，再检查代表文档是否支持。若合并两簇，计数变成$(4,2,2)$，需按新簇定义重新计算平均长度及代表词，不能只拼接旧名称。

新文档沿固定编码器、UMAP变换和所用HDBSCAN预测接口归属，可能继续成为离群项；某些配置没有可靠的外样本预测，需要另设分配器并独立验收。簇成员得分反映所用密度结构，不是pLSA/LDA的词生成主题比例。将混合文档硬分到A，不意味着其“场地”部分不存在；生成的簇名也不能作为每个成员的逐句事实。

\subsubsection{主题数选择、解释与评价}
\label{sec:nlp-classification-topic-evaluation}

主题词看起来合理，只说明通过了一次人工印象检查。连贯性指标考察高权重词在参考语料中是否共同出现；主题词多样性衡量不同主题是否反复使用同一组词；跨运行稳定性先对齐主题，再比较词分布或文档归属。词语侵入测试要求读者找出混入主题词列表的异类，文档侵入测试要求找出不属于一组材料的文档，补充自动指标无法代替的解释性证据。

例如在100个固定窗口中，“讲座”出现20次，“报名”30次，两者共同出现10次。归一化点互信息为
\[
 \operatorname{NPMI}(u,v)=
 \frac{\log[p(u,v)/(p(u)p(v))]}{-\log p(u,v)}
 =\frac{\log(0.1/(0.2\times0.3))}{-\log0.1}
 \approx0.2218.
\]
窗口长度、参考语料和平滑必须固定。共现为零时公式不能直接取对数，应按评价协议平滑或使用边界约定；词对评分再如何聚合成主题评分也须说明。即使词经常同现，也可能只是模板页脚，而不是读者需要的主题。

生成式模型还可以比较留出词的预测似然或每词困惑度，但词表、分词及推断条件应一致。文档补全评价可用一部分词推断新文档，再对另一部分词计分；用全部词拟合$\theta$后又评价这些词，不是同一留出协议。BERTopic等聚类路线没有天然的LDA式生成似然，不应强行给它一个同名困惑度。

选择主题数时，固定语料和预处理，比较若干$K$下的重复主题、主题词多样性、稳定性和人工用途；密度聚类则记录最终簇数和离群率怎样随参数变化。数量过少会合并讲座和场地，过多可能把同一活动的模板拆成多个簇。短文本、罕见缩写、多义词和混合通知应单列检查。无标签主题适合探索目录，若后续必须与业务栏目一一对应，还需要独立映射及监督验收。

\subsection{文档聚类}
\label{sec:nlp-classification-clustering}

文档聚类将表示相近的材料放在一起，但“相近”由词表、编码器、归一化和距离共同决定。TF-IDF可能强调术语，SBERT可能强调整体语义，日期与活动编号是否主导分组取决于当前需求。主题聚类允许不同日期的同类活动放在一起，版本聚类却必须先区分活动身份，不能共用一个未经说明的距离。

本节把四份通知的二维表示固定为$A=(1,0),B=(0.8,0.2),C=(0,1),D=(0.2,0.8)$，另以$E=(0.5,0.5)$代表混合材料。向量均为教学构造；返回文档时要保留ID，选靠近中心或代表密集区域的真实文档解释，不能把一个抽象均值当作原文。

\subsubsection{K-means与球面K-means文档聚类}
\label{sec:nlp-classification-kmeans}

普通K-means以平方欧氏距离分配文档，再取簇内均值更新中心。对$A,B,C,D$取两个初始中心$(1,0),(0,1)$，B到它们的平方距离为0.08和1.28，因此分到第一簇；D对称地分到第二簇。更新中心得到$(0.9,0.1)$和$(0.1,0.9)$，四点总平方误差从0.16降至0.08。重新分配不变后，均值也不变，本例达到稳定。通用优化及局部最优问题见第~\ref{chap:clustering}章。

球面K-means先将非零文档向量归一化为单位向量，以余弦最大分配，中心更新为簇内向量和的单位化。B变为$(0.9701,0.2425)$；A与B之和再归一化，得到约$(0.9925,0.1222)$，不是普通均值$(0.9,0.1)$。对单位文档与单位中心有$\|\bm x-\bm c\|_2^2=2-2\bm x^\top\bm c$，因此最大余弦与最小该距离等价；中心不归一化时不能直接套用。

新文档可以按最近中心分配，但E到两个普通中心同距，应使用事先规定的并列规则或保留低确定性，不能硬称它只有一个主题。改变TF-IDF、LSA或SBERT表示，以及随机初始化、$K$和归一化，均可能改变簇。停止条件还应包含最大迭代数、中心移动或目标变化；空簇需按预设规则重启。

\subsubsection{凝聚层次聚类}
\label{sec:nlp-classification-hierarchical}

凝聚层次聚类从每篇一个簇开始，反复合并最近两簇，保留每次合并的距离形成树。单链接取跨簇最近点距离，全链接取最远点距离，平均链接取全部跨簇点对距离的平均；它们回答的是不同的“组与组多近”，不能只更换名称而沿用同一结果。

对四份通知，$A,B$距离为$\sqrt{0.08}\approx0.2828$，$C,D$相同，先形成两簇，距离并列按固定ID顺序处理。两簇间单链接距离为$B,D$的$\sqrt{0.72}\approx0.8485$，全链接为$A,C$的$\sqrt2\approx1.4142$，平均链接为约1.1314。若在距离1.0截断，单链接已合成一簇，另两种仍保留两簇。加入桥接文档E时，单链接更容易通过局部近邻串成大组。

解释树时选择代表文档并说明截断标准。算法树表示合并顺序，不自动意味着“活动通知—讲座—机器学习”这样的业务包含层级。经典凝聚算法没有固定参数的天然新文档后验；重新聚类或另设最近簇分配器是不同使用方案，应记录会不会改变已有目录。

\subsubsection{DBSCAN与HDBSCAN密度聚类}
\label{sec:nlp-classification-density}

密度方法适合保留“不属于任何已发现组”的材料。DBSCAN以半径$\varepsilon$和最少邻域点数$\mathrm{MinPts}$定义核心点，再沿核心点的密度可达关系恢复簇，边界点与噪声按规则处理。对四点加E，取欧氏距离、$\varepsilon=0.3,\mathrm{MinPts}=2$且邻域计入自身，A/B和C/D各形成一簇；E到最近的B、D均约0.4243，只有自身，成为噪声。

把$\varepsilon$增至0.45，E连到B和D，核心点连接可把两组并成一簇。这是密度连通变化，不是文本突然变成同一主题。边界点同时邻接不同簇时，实现的遍历与归属约定也可能影响结果，参数与输入顺序应记录。

HDBSCAN利用核心距离与互达距离建立密度层次，再按稳定性等规则选择簇，减少对单一半径的依赖。以$\mathrm{MinPts}=2$的上述距离看，A的核心距离约0.2828，E约0.4243；互达距离取两端核心距离与点间距离的最大值，E连接各组的边因而处在较稀疏层级。完整层次构造与簇选择回用聚类章节，不把HDBSCAN简化成“自动选一个DBSCAN半径”。

小簇、稀疏领域和混合材料可能被作为离群项，不一定是坏数据。新文档能否近似归属依实现是否保留预测信息；重新拟合则可能改变原簇。评价时同时报告覆盖率和噪声比例，轮廓系数可在非噪声子集上算，但必须注明排除了谁，不能隐藏大量难文档后只报告较好得分。

\subsubsection{高斯混合与文档软分组}
\label{sec:nlp-classification-gmm}

若要保留多个组的可能性，可在连续文档表示上拟合高斯混合
$p(\bm x)=\sum_k\pi_k\mathcal N(\bm x;\bm\mu_k,\bm\Sigma_k)$。
EM以成分后验为责任度，更新混合权重、均值和协方差；协方差取全矩阵、对角或共享形式会改变分组边界，高维小样本下常需正则化。它建模的是文档向量分布，不是词逐个来自主题的过程。

取两成分等先验，均值为$(0.9,0.1)$与$(0.1,0.9)$，共享协方差$0.08\bm I$。E到两个中心平方距离均为0.32，后验各0.5。新通知$F=(0.7,0.3)$的两距离为0.08和0.72，密度比为$\exp[(0.72-0.08)/(2\times0.08)]=e^4$，第一成分后验为$\sigma(4)\approx0.9820$。可保留这个软向量，或按最大后验分组；接受阈值需另行验证。

E内容混合与后验0.5可以同时出现，但二者不是逻辑等价：一篇离两簇都很远的文档也可能后验各半。因此未知文档还应检查绝对密度或独立拒绝条件。这里的0.5是“一个文档向量属于哪个成分”的后验，不是“半数词由哪个主题生成”的比例。

各路线共同使用轮廓系数、跨运行稳定性和代表文档核验；有参考分组时可补充ARI、NMI及业务一致性。轮廓系数$s_i=(b_i-a_i)/\max(a_i,b_i)$比较同簇平均距离$a_i$与最近其他簇平均距离$b_i$，单点簇、只剩一簇及噪声的处理必须约定。好看的簇名与较高内部指标都不能替代未知文档分配的实际验证。

\subsection{关键词与关键短语组织}
\label{sec:nlp-classification-keyphrases}

文档分组后，用户仍需要简短入口找到材料。\term{关键词抽取}选择原文中的词或跨度，\term{关键短语生成}则允许产生原文未出现的概括。单篇关键词服务当前材料，主题代表词解释潜在分布，簇代表词解释文档组；相同词语在不同接口下的得分并不相同。

一个可检查的流程需要明确候选从何而来、如何计分、怎样去重，以及如何恢复到原文或合法生成短语。候选阶段遗漏“报名截止时间”，后面的排序无法补回；把“机器、学习、报名”排在前面，也不等于已经生成了语法完整的短语或全文摘要。

\subsubsection{TF-IDF：词频与跨文档区分度}
\label{sec:nlp-classification-keyword-tfidf}

TF-IDF用单篇词频与参考语料的文档频率区分常见词和局部重要词。沿用$\operatorname{idf}(v)=\log((N+1)/(\operatorname{df}_v+1))+1$，若参考集有4篇，“读书会”只出现在1篇、在当前通知出现2次，权重为$2(\log(5/2)+1)\approx3.8326$；“通知”在4篇都出现，当前出现3次，权重为3。前者虽然次数较少，仍因跨文档区分度更高排在前。

参考语料必须说明来自哪里，不能声称仅看当前单篇就已经知道跨文档IDF。停用词过滤可减少虚词，过度低频过滤却会删掉活动编号和新领域术语。长度归一化在单篇内部可能不改变排序，但会影响跨文档比较。

短语可以由词性模式、连续$n$-gram或已知术语表生成，再按整短语文档频率计分，或按规定聚合组成词权重。两者不等价：高分的“机器”和“报名”若不相邻，就不能拼成一条抽取短语。应保留候选原始边界、长度规则和出现次数，评价完整短语而不是只命中其中一个词。

\subsubsection{RAKE：候选切分与词共现度}
\label{sec:nlp-classification-rake}

RAKE利用停用词和分隔符把文档切成候选短语，再由短语内的词共现给候选排序。词频$f(w)$计出现次数，共现度$\deg(w)$在这里采用包含自身的口径；一次长度$r$的候选对其中每个词贡献$r$。词分数为$\deg(w)/f(w)$，短语分数为其词分数之和。若实现先只数其他词，需再加自身频次，不能混淆两种度的定义\scite{nlp-rose2010rake}

从教学通知中固定候选“机器／学习／读书会”“报名／信息”，再加入同一文档另一处“机器／学习”。“机器、学习”各出现2次、度各为$3+2=5$，词分数各2.5；“读书会”出现1次、度为3，分数3；“报名、信息”各出现1次、度为2，分数各2。三个候选得分依次为8、4、5，因此排序为完整活动短语、较短的“机器学习”、报名信息。

输出时去掉精确重复，若导航预算不希望同时保留包含关系候选，可保留完整短语并删去“机器学习”，但必须声明这个后处理规则。原方法也讨论反复相邻候选的组合；不是把所有停用词两边的词无条件重新拼回。中文分词与停用词表决定“机器学习”是一词还是两词，改变切分就要重算，不能移用英语资源却把差异归因于语义能力。

\subsubsection{YAKE：单篇文本的局部统计}
\label{sec:nlp-classification-yake}

没有参考语料时，YAKE使用当前文档中的大小写、出现位置、频率、分布范围和局部上下文形成分数。与TF-IDF相反，YAKE的候选分数通常越小越优；按降序输出会颠倒结果。大小写特征主要适用于有字母大小写的语言，中文文本不能凭空获得这种信号，需记录分词与语言配置\scite{nlp-campos2020yake}

以公开实现中的组合说明计算。设$C$为大写频次与非句首首字母大写频次中的较大者除以$1+\log f$，$F=f/(\overline f+s_f)$为频率特征，$\overline f,s_f$是非停用词词频的均值和标准差，$D$为出现句数占比，$P=\log\log(3+\text{句子索引中位数})$为位置特征，句索引从零开始。令
$R=1+(L+R_{\!c})f/f_{\max}$，其中$L,R_{\!c}$分别是左、右上下文的不同邻词数与邻接总次数之比，无邻词时取零。单词分数组合为
\[
 H(w)=\frac{P R}{C+F/R+D/R}.
\]
较早出现使$P$较小；频繁且分布多句可增大分母；邻接关系更分散则会影响$R$。这些是设计的统计启发，不是关键词真值的概率。

给一个词的完整教学统计：$f=2,f_{\max}=4,\overline f=2,s_f=1$，在4句中的第0、1句出现，且$L=R_{\!c}=0.5$，没有大小写信号。于是$C=0,F=2/3,D=1/2,P=\log\log3.5,R=1.5$，得到$H\approx0.4346$。若与另一词组成不含停用词的候选，另一词$H=0.3$，短语只出现1次，则按
\[
 H(c)=\frac{\prod_{w\in c}H(w)}
             {f(c)\,[1+\sum_{w\in c}H(w)]}
\]
得分约0.0752。含内部停用词的候选还需按所用实现作邻接修正，不能直接套这个无停用词例。

实际通知无需外部IDF就能产生局部候选，但标题重复、模板开头和领域缩写可能被过度奖励。最大$n$-gram长度、首尾停用词过滤及近似字符串去重阈值都影响输出。应同时查看正文位置和被合并的候选，防止将“报名时间”与“活动时间”因字面相近误合并。

\subsubsection{TextRank：共现图与关键词排序}
\label{sec:nlp-classification-textrank}

TextRank将候选词作为节点，在固定窗口中共现的词之间连边，以局部连接传播重要性。先规定分词、词性筛选、窗口和边权，再采用PageRank迭代。这里将分数归一化为和为1，对$M$个节点使用
\[
 r_i^{(t+1)}=\frac{1-d}{M}
 +d\sum_{j:\,w_{j,i}>0}
 \frac{w_{j,i}}{\sum_k w_{j,k}}r_j^{(t)}.
\]
阻尼$d\in(0,1)$，无出边节点按均匀分布向全部节点转移；采用未归一化常数$(1-d)$的写法会改变尺度，但不应混在同一数值例里\scite{nlp-mihalcea2004textrank}

取“机器、学习、读书会”三个节点，只有无向边“机器—学习”和“学习—读书会”，权重均1。初值各$1/3$，$d=0.85$时，一轮后两端各为$0.05+0.85/6\approx0.1917$，中间为$0.05+0.85(2/3)\approx0.6167$，和仍为1。继续迭代至相邻向量$L_1$差小于容差，才得到稳定排名；不能把一轮结果当成收敛值。

选定高分词后，根据原文连续出现和候选模式恢复“机器学习读书会”，而不是按分数把不相邻词随意连接。窗口过大容易让通用词连接许多对象，词性过滤也可能删去关键动词。第~\ref{chap:nlp-text-generation}章的句子选择可以复用图排序思路，但节点变成句子，输出也变成摘要句集合，不是同一个关键词任务。

\subsubsection{EmbedRank：语义表示与MMR选择}
\label{sec:nlp-classification-embedrank}

EmbedRank先用名词和形容词等规则取得候选短语，将候选及文档编码，以余弦相似度排序。原工作考察Doc2Vec和Sent2Vec等句级表示，并按候选词性过滤文档表示；其增加多样性选择的版本称为EmbedRank++，用最大边际相关性（MMR）平衡相关与重复。它不是训练一个有监督关键词分类头；\citeyear{nlp-bennani2018embedrank}年原始论文的完整出处见本章文献说明。

为简化手算，下面直接用原始余弦，区别于原文对相似度的规范化处理。已有选中集合$S$时，从未选候选中最大化
\[
 \operatorname{MMR}(c)=\lambda s(c,x)
 -(1-\lambda)\max_{c'\in S}s(c,c').
\]
首个候选按文档相似度选择。取原文候选A“机器学习读书会”、B“读书会”、C“报名信息”，与文档的相似度分别为0.90、0.88、0.75；A/B相似度为0.95，A/C为0.40。先选A，$\lambda=0.7$时，B得$0.7(0.88)-0.3(0.95)=0.331$，C得$0.7(0.75)-0.3(0.40)=0.405$，第二项选C。若$\lambda=0.9$，B为0.697、C为0.635，选择又变回B。

这个变化反映导航需求对重复的容忍度，不是哪个$\lambda$天然正确。“报名信息”虽能提示报名事项，却没有保留真正重要的截止时刻；应核查所选候选是否覆盖关键事项。抽取式候选须能定位到原文，若另行生成原文没有的“报名截止时间”作为概括名称，应标明生成或概括候选，不能报告为原文抽取命中。

\subsubsection{序列标注与跨度式关键短语抽取}
\label{sec:nlp-classification-keyphrase-spans}

有参考关键短语时，可以将其原文出现位置转换成BIO标签，复用编码器和CRF；也可以枚举有限长度跨度，直接预测边界和关键性。标注规范应先决定：同一短语出现多次是否全部标注，嵌套短语是否都保留，参考中未出现于原文的概括词是否排除出抽取训练。通用标注和跨度模型分别见第~\ref{chap:nlp-text-analysis}章与第~\ref{chap:nlp-information-extraction}章。

对“机器／学习／读书会／报名／截止”，若参考跨度为前3词及后2词，BIO序列是$(B,I,I,B,I)$。独立词分类若把第4词误为$O$、第5词为$I$，恢复时就出现非法开头；CRF可以用转移约束禁止$O\to I$，或用明确后处理修正，但应重新评价边界。跨度模型则可给$(1,3)$和$(4,5)$分数0.92、0.86，在验证阈值0.8下恢复两段，并映射回字符区间。

若同时需要“机器学习”和“机器学习读书会”，一套平坦BIO标签不能表达嵌套，需多层标签或跨度集合；使用不重叠解码则会主动舍弃其中一个。分别报告候选覆盖、完整边界F1和去重后的短语F1，才能区分未找到位置、边界不全和选错短语。仅给定正确跨度后的分类准确率，不代表端到端抽取质量。

\subsubsection{CopyRNN：生成与复制式关键短语预测}
\label{sec:nlp-classification-copyrnn}

抽取器不能输出原文未出现的概括短语。CopyRNN以“文档—单条关键短语”为训练对，使用双向GRU编码与带注意力的GRU解码；同一文档有多个参考短语时形成多个训练对。复制分支让模型可以输出原文中的低频词或未登录词，生成分支则可以组合词表中的概括词。其\citeyear{nlp-meng2017copyrnn}年原始论文的完整出处见本章文献说明。

在一步解码中，生成词表项有logit $g_t(v)$，原文位置有复制logit $c_t(j)$。原方法把两类未归一化质量放进同一个分母：
\[
 \begin{aligned}
   Z_t&=\sum_{v\in V}\exp g_t(v)+\sum_{j=1}^{|x|}\exp c_t(j),\\
   p_t(w)&=\frac{
     \mathbf1[w\in V]\exp g_t(w)
     +\sum_{j:x_j=w}\exp c_t(j)}{Z_t}.
 \end{aligned}
\]
一个词在原文多处出现时，复制质量相加；同时在词表和原文出现时，两路都可贡献。这里不是两个分别归一化的分布直接相加，也不是另引一个未在该配置中定义的混合门。

例如一步中词表仅展示“报名、信息、结束符”三项，生成质量分别为$(2,1,1)$，其余生成项质量在这个教学小词表中为零；原文三个位置为“海棠楼、报名、报名”，复制质量为$(1,1,2)$。共同分母为8，故$p(\text{报名})=(2+1+2)/8=0.625$，$p(\text{海棠楼})=1/8$，信息与结束符也各$1/8$。虽然“海棠楼”不在生成词表，它仍可通过原文位置输出。合并同词的位置质量以后，词级分布才用于扩展候选。

训练目标是参考短语及结束符的负对数似然，前一步使用参考词作条件。若“报名／信息／结束符”三步的条件概率为$(0.625,0.5,0.8)$，序列概率为0.25，损失为$-\log0.25\approx1.3863$；每步都要用该步状态重新计算$Z_t$。这不是三条独立关键词的概率。

运行时用开始符启动，自回归地扩展词，并以结束符停止一条短语。束搜索保留若干高累计分数前缀，再从已完成序列取得多个候选。只看两条教学路径，“报名／结束符”若概率为$0.625\times0.3=0.1875$，“报名／信息／结束符”为0.25，按未归一化序列概率后者更高；实际束宽、长度惩罚和最大长度会影响能否保留它们，必须连同输出记录。去掉结束符后还需规范化、去重，核对短语支持范围。

这种单短语训练与把整组短语按某个次序串起来生成不同：后者需要集合到序列的顺序和分隔符规范，并会引入短语间条件依赖。评价应分开原文中出现和未出现的参考短语，声明精确匹配或词形规范化；匹配参考、得到原文支持和适合导航是三个条件。能复制地点不证明选它最有用，能生成流畅概括也不证明原文支持该概括。

\section*{文献说明}

以下列出文本生成参与判定、ESIM、BERTopic以及两种关键短语方法的原始资料。
类别定义、模型配置与输出含义分别依据对应文献。
\begin{fullwidth}
  \small
  \noindent\fullcite{nlp-li2026creator}\par\medskip
  \noindent\fullcite{nlp-chen2017esim}\par\medskip
  \noindent\fullcite{nlp-grootendorst2022bertopic}\par\medskip
  \noindent\fullcite{nlp-bennani2018embedrank}\par\medskip
  \noindent\fullcite{nlp-meng2017copyrnn}
\end{fullwidth}

\section{本章小结}
\label{sec:nlp-classification-summary}

一份通知能被贴标签、与其他文字建立关系，也能参与集合组织，三个层次回答的是不同问题。属性判定先规定内容、情感、立场或来源标签，单标签与多标签采用不同输出约束，训练损失之外还要冻结验证得到的阈值。对人机共同编辑或未知来源，统计分数尤其不能升级成身份的确定证据。

句对判断需要把关系写清。对称相似度不能替代有方向的蕴含，高词汇重合不能保证复述，主题相关也不保证能回答查询。双编码、交叉编码和后期交互提供不同计算接口，候选召回、关系判定与原文恢复分别可能失败。对齐要保留未配对及一对多，近重复还要核对日期、实体和版本来源。

没有预设类别时，可以从词共现、潜在表示和文档距离寻找结构。主题比例、非负系数、潜在坐标和聚类成员分数各有模型内含义，不应统称为主题概率。关键词则将材料组织成可浏览入口，候选、排序、去重和事实支持都须核查。下一章转向生成与转换：当输出不再只是标签、配对或原文片段，而是新写成的文字时，还需要控制它保留了什么、改写了什么，以及是否引入原文没有的信息。
